% language=en

\usemodule[article-basic]
\usemodule[abbreviations-mixed]
\usemodule[scite]

\setuplayout
  [article]
  [backspace={2*\measure{article:margin}},
   topspace={1.5*\measure{article:margin}},
   bottomspace={1.5*\measure{article:margin}}]

\setupalign[mathbookpasses]

\definedescription
  [keyword]
  [headcolor=darkred,
   headstyle=,
   before=,
   after=,
   alternative=serried,
   width=fit]

\startsetups mixedsetups
  % \switchtobodyfont[9pt]
    \setupwhitespace[none]
\stopsetups

\setupmixedcolumns
  [n=2,
%    splitmethod=none,
%    grid=tolerant,
   balance=yes,
   setups=mixedsetups,
   splitmethod=fixed,
   grid=yes,
   internalgrid=line]

\setupcombination
  [distance=2ts]

\setupcaptions
  [width=0.8tw]

\setuphead
  [section]
  [page=yes]

\setupinteraction
  [state=start,
   focus=standard,
   color=,
   contrastcolor=]

\defineframed
  [TightBox]
  [offset=overlay]

\setupregister
  [n=3,
   color=darkred]

\setupregister
  [index]
  [pagestyle=]

\defineregister
  [keywordindex]

\setupregister
  [keywordindex]
  [textstyle=\tt,
   pagestyle=]

\starttext

% We want proper coloring so we wrap in a TEX page instead, scite side effect
% would color "color" in the setup somewhat strange .. maybe a todo.

\startbuffer[coverpage]
\startMPcode
StartPage ;

picture q ; q := image (

    lmt_scene_start [
        width           = 18cm,
        crop = true,
        projection      = "perspective",
        eye             = (2.7,-4.0,2.4),
        light           = (0,0,-1),
        intensity       = 1,
    ] ;

    lmt_scene_material [
        name    = "sphere",
        diffuse = (1,1,1),
    ] ;

    lmt_scene_parametric [
        id       = "sphere",
        x        = "sin(u)*cos(v)",
        y        = "sin(u)*sin(v)",
        z        = "cos(u)",
        umin     = 0,
        umax     = 2pi,
        vmin     = 0,
        vmax     = 2*pi,
        nu       = 5,
        nv       = 5,
        material = "sphere",
        normal   = "flat",
    ] ;

    lmt_scene_stop ;
) ;

q := q shifted - center q
       shifted (0.5PaperWidth,0.5PaperHeight) ;

draw q ;

picture t[] ;
t[0] := textext("\sansbold lmt_scene") ;
t[1] := textext("\sansbold Hans Hagen & Mikael Sundqvist") ;

t[0] := t[0] xsized .6 PaperWidth ;
t[1] := t[1] xsized .6 PaperWidth ;

t[0] := t[0]
    shifted - ulcorner t[0]
    shifted (1cm,-1cm)
    shifted (0*PaperWidth,1*PaperHeight) ;

t[1] := t[1]
    shifted - lrcorner t[1]
    shifted (-1cm,1cm)
    shifted (1*PaperWidth,0*PaperHeight) ;

draw t[0] withcolor red ;
draw t[1] withcolor red ;

StopPage ;
\stopMPcode
\stopbuffer

\startTEXpage[pagestate=start,background=color,backgroundcolor=black]
    \getbuffer[coverpage]
\stopTEXpage

\completecontent

\startsection
  [title=Introduction]

\METAFUN\ and \METAPOST\ are mainly two-dimensional drawing tools, but many
figures also need a three-dimensional view. The \type {lmt_scene} commands give
you one place to define the geometry, camera, lighting, and annotations for that
view. In this document we collected what is currently available. Since the project
started out as an AI experiment, we also let Chat GPT (Luna) generate parts of the
content of this manual.

The interface is scene-based. With a scene you can:

\startitemize
  \startitem
    choose the camera, light, colors, output size, and other shared settings;
  \stopitem
  \startitem
    draw several kinds of geometry: surfaces given by a
    function, \m {z = f(x,y)}, parametric surfaces given by \m{(x,y,z) = (f _ 1(u,v), f _
    2(u,v), f _ 3(u,v))}, implicit plots described by \m {F(x,y,z) = 0}, parametric curves \m
    {(x,y,z) = (f _ 1(t), f _ 2(t), f _ 3(t))}, finite plane patches
    given by a center point, a normal direction, and two sizes, and imported
    triangle meshes from external files.
  \stopitem
  \startitem
    add reference geometry such as grids, axes, ticks, and annotations; and
  \stopitem
  \startitem
    combine the objects with common camera and depth handling.
  \stopitem

\stopitemize

A figure starts as a description of its objects, camera, materials, and output
settings. Surfaces are sampled over chosen domains, while imported or generated
geometry is added as meshes. All of these objects share one camera and one
depth-aware renderer, so they appear in the same picture. After the scene is
rendered, projection helpers can place ordinary \METAPOST\ labels and
annotations.

Lighting is simple: a directional light works with the material color, ambient
level, highlight, and transparency. These few controls are enough to show the
shape of an object without making the setup hard to follow. Stippling is a raster
postprocessor and is covered near the end of the manual.

The labels in the table of contents are hints about when to read a section: \emph
{optional} sections can wait, \emph {advanced} sections assume the basic scene
workflow, and \emph {experimental} sections describe interfaces that may change.
None of them is needed for a first scene.

\startsubsubject[title=Roadmap]

This route gives a useful starting order:

% normally \in but then we need \dontleavehmode ... added now
% we could use \about here

\startitemize
  \startitem
    \goto {Quick start} [quick-start] shows the smallest complete file and the
    three-step scene lifecycle.
  \stopitem
  \startitem
    \goto {Scene basics} [scenes] explains the shared camera, output, lighting,
    and cropping settings; \goto {Defaults} [defaults] collects the principal
    defaults in one place.
  \stopitem
  \startitem
    \goto {Materials} [materials] explains how to control color, lighting,
    transparency, and line-like helpers.
  \stopitem
  \startitem
    \goto {Surface graphs} [surface-graphs] is the usual starting point when the
    geometry is given by \m {z=f(x,y)}.
  \stopitem
  \startitem
    Continue with \goto {parametric surfaces} [parametric-surfaces],
    \goto {implicit surfaces} [implicit-surfaces],
    \goto {curves} [parametric-curves],
    and \goto {intersections} [intersections] as the geometry requires. The
    \goto {Reference geometry} [reference-geometry] section covers
    \goto {planes} [planes], \goto {segments} [segments], \goto {grids} [grids],
    \goto {axes} [axes], and \goto {ticks} [ticks].
  \stopitem
  \startitem
    \goto {Rendering pipeline} [rendering-pipeline] explains why mesh counts,
    raster resolution, and supersampling affect different stages of a render.
    The \goto {Putting it all together} [complete-scene] example then combines
    the main scene features in one file.
  \stopitem
  \startitem
    Leave \goto {caching} [scene-cache],
    \goto {annotations} [annotation-helpers], \goto {stippling} [stippling],
    \goto {special overlays} [special-overlays],
    \goto {low-level \LUA\ helpers} [lua-helpers], and
    \goto {mesh inspection} [mesh-inspection] until the basic scene workflow is
    familiar. External and internal meshes are introduced in the object table
    when those forms of geometry are needed.
  \stopitem
\stopitemize

Command sections explain the purpose of an interface, show one or more runnable
examples, and collect the supported keywords nearby.

\stopsubsubject

\startsubsubject[reference=quick-start,title=Quick start]

To try a scene, save the complete file below as \type {hello3d.tex} and run
\type {context hello3d.tex} in the same directory. You need a \CONTEXT\
installation that includes the \LUAMETATEX\ scene commands. The \type {offset}
leaves a little space around the rendered page. Later examples use buffers so
that their source and result can be shown together.

\startbuffer[first-picture]
lmt_scene_start [
  width  = 7cm,
  height = 5cm,
  crop   = true,
] ;

lmt_scene_surface [
  code = "sin(x)*cos(y)",
  xmin = -pi,
  xmax =  pi,
  ymin = -pi,
  ymax =  pi,
] ;

lmt_scene_stop ;
\stopbuffer

\startbuffer[first-picture-file]
\starttext

\startMPpage[offset=1cm]
lmt_scene_start [
  width  = 7cm,
  height = 5cm,
  crop   = true,
] ;

lmt_scene_surface [
  code = "sin(x)*cos(y)",
  xmin = -pi,
  xmax =  pi,
  ymin = -pi,
  ymax =  pi,
] ;

lmt_scene_stop ;
\stopMPpage

\stoptext
\stopbuffer

\typebuffer[first-picture-file][option=TEX]

\startplacefloat
  [figure]
  [title={A first scene: the graph of \m {z=\sin(x)\cos(y)}.}]
  \processMPbuffer[first-picture]
\stopplacefloat

This manual describes the scene module generated on 2026-05-01 or newer.
Older installations may report unknown \type {lmt_scene_*} commands or
keywords in the log; update \CONTEXT\ when that happens.

The scene is built in three steps: start it, add one or more objects, and stop
it. The \type {code} value is a quoted expression because it is evaluated while
the surface is being generated. The names\index {expressions} \type {x}, \type
{y}, and \type {z} refer to coordinates; the available operators, functions, and
constants are collected in the \goto {Expression language} [expressions]
subsection below.

There are a few conventions worth remembering:

\startitemize
  \startitem
    A triplet such as \type {(1,0,0)} is a 3D point or direction. Its meaning
    depends on the keyword that receives it.
  \stopitem
  \startitem
    Material and background colors are \RGB\ triplets with components between \type
    {0} and \type {1}; for example, \type {(1,0,0)} is red. Line helpers use the
    same color values through their \type {color} keyword; a \METAPOST\ color
    name such as \type {red} is accepted there as well. Even better, use colors
    that are shared with the \TEX\ end, like \type {"red"}.
  \stopitem
  \startitem
    Give an object an \type {id} when another object or a query must refer to it.
    An intersection can only refer to objects that have already been added.
  \stopitem
  \startitem
    Start with the default camera and material. Add \type {eye}, \type {target},
    and a named material only when the first picture needs adjustment.
  \stopitem
\stopitemize

Keep the first example small. Once it works, change one setting at a time. Add a
named material when the default appearance is not enough, adjust \type {eye} or
\type {target} when the view is wrong, and increase the mesh counts or raster
resolution only when the result needs more detail. The later sections build on
this same workflow.

\stopsubsubject

\startsubsubject[title=Build on the first scene]

Now give the surface a name and a material. The geometry stays the same; only the
view and appearance change:

\starttyping[option=MP]
lmt_scene_start [
  width  = 7cm,
  height = 5cm,
  crop   = true,
  eye    = (3,-4,2.5),
  target = (0,0,0),
] ;

lmt_scene_material [
  name    = "wave",
  diffuse = (.15,.55,.90),
  ambient = .25,
] ;

lmt_scene_surface [
  id       = "wave",
  code     = "sin(x)*cos(y)",
  xmin     = -pi,
  xmax     =  pi,
  ymin     = -pi,
  ymax     =  pi,
  material = "wave",
] ;

lmt_scene_stop ;
\stoptyping

The \goto {Materials} [materials] section explains the appearance settings. The
surface sections show how to replace this graph with a parametric or implicit
object. The sections on planes, intersections, and reference geometry then show
how to turn the picture into an annotated figure.

\stopsubsubject

\startsubsubject[title=Common adjustments]

When the picture needs tuning, change one group of settings at a time:

\startitemize
  \startitem
    If the surface looks faceted or misses a small feature, increase the object
    mesh counts such as \type {nx}, \type {ny}, \type {nu}, or \type {nv}.
    For an implicit surface, increase all three of \type {nx}, \type {ny}, and
    \type {nz}.
  \stopitem
  \startitem
    \index {mesh counts}%
    \index {raster resolution}%
    \index {supersampling}%
    If edges look jagged, increase \type {resolution} or try
    \type {supersample = 2} or \type {4}. These settings affect the raster,
    not the geometric mesh.
  \stopitem
  \startitem
    If the object is seen from the wrong side, adjust \type {eye} and
    \type {target}. Use \type {crop = false} while comparing camera settings.
  \stopitem
  \startitem
    If the shape looks flat, try another \type {light} direction, adjust the
    material's \type {ambient} value, or choose a smoother normal mode.
  \stopitem
  \startitem
    If a plane should sit over the rest of the scene, give it a low
    \type {opacity}; if a line should remain visible, use a material with
    \type {ontop = true}.
  \stopitem
\stopitemize

\stopsubsubject

\startsubsubject[reference=expressions,title=Expression language]

\index {expressions}%
The quoted scene expressions (or course) use the \LUA\ expression language.
Arithmetic uses \type {+}, \type {-}, \type {*}, \type {/}, the remainder
operator, and \type {^} for powers; comparisons use \type {==}, \type {~=}, \type
{<}, \type {<=}, \type {>}, and \type {>=}; and the logical operators are \type
{and}, \type {or}, and \type {not}. The remainder operator is written as follows
in the expression itself:

\starttyping
x % y
\stoptyping

Parentheses can be used wherever they make the order of evaluation clearer.

Angles are in radians. The constant \type {pi} is available, so a full turn is
normally written \type {2*pi}; there is no degree conversion in the scene
expression syntax.

The expression environment exposes \LUAMETATEX's extended math table as
\type {math} and also makes its names available at the top level. Useful names
include \type {sin}, \type {cos}, \type {tan}, \type {asin}, \type {acos},
\type {atan}, \type {atan2}, \type {sinh}, \type {cosh}, \type {tanh},
\type {exp}, \type {log}, \type {log10}, \type {sqrt}, \type {cbrt},
\type {floor}, \type {ceil}, \type {round}, and \type {trunc}. For absolute
values and extrema, use \type {abs}/\type {fabs}, \type {min}/\type {fmin},
and \type {max}/\type {fmax}; the \type {f...} spellings are the portable
names in the extended \type {xmath} table. Related helpers include
\type {hypot}, \type {pow}, \type {deg}, \type {rad}, \type {dotproduct},
\type {crossproduct}, and \type {normalize}.

When you use functions, be aware that the default environment is \type {xmath}
but we also register \type {MP} (the use namespace), \typ {documentdata}, \type
{userdata}, \type {moduledata} and \type {globaldata}. By making \type {xmath}
default performance remain good.

The variables depend on the object command. Surface graph expressions use
\type {x} and \type {y}; parametric-surface coordinates and tangent expressions
use \type {u} and \type {v}; implicit expressions and derivatives use
\type {x}, \type {y}, and \type {z}; and curve coordinates use \type {t}.
The \LUA\ \type {pfunction} and \type {nfunction} forms receive \type {x},
\type {y}, and \type {z} as function arguments. The \type {z} value of a graph
is produced by its expression and is not an independent input variable.

Expression loading and evaluation diagnostics are written to the job's \type
{.log} file. A failed expression is reported with its expression name; inspect
the log before changing the sampling grid or camera.

\stopsubsubject

\stopsection

\startsection
  [reference=scenes,title=Scene basics]

A scene contains one 3D picture. It holds the objects and the shared settings
that control them: output size, camera, lighting, background, and quality.
Surfaces, curves, grids, axes, and other helpers can all share the same scene.

Most scenes have this shape:

\starttyping[option=MP]
lmt_scene_start [ ... ] ; % open the scene
  % here we add our 3D graphics
lmt_scene_render ; % optional: render before annotations
  % here we add extra material
lmt_scene_stop ; % finish the scene
\stoptyping

\type {lmt_scene_start} opens a scene and records its shared settings. The object
commands between start and stop add geometry. Most figures need only
\type {lmt_scene_stop}. Use \type {lmt_scene_render} when later \METAPOST\ code
needs the rendered size, projected positions, or visibility information before
the scene is closed.
\index {bytemap}%
Stopping the scene places the finished bytemap in the
current picture.

Here is a small initialization:

\starttyping[option=MP]
lmt_scene_start [
  width       = 5cm,
  projection  = "perspective",
  eye         = (3.0,-4.2,2.2),
  target      = (0,0,0),
  crop        = true,
  light       = (-1,-1,-2),
] ;
\stoptyping

This scene is five centimeters wide and uses perspective projection. The camera
is at \type {(3.0,-4.2,2.2)}, and the light comes from \type {(-1,-1,-2)}. The
other scene options fall into four groups: output size and quality, camera and
view, lighting and background, and optional postprocessing.

\index {cameras}%
Camera and output settings control different things. The camera decides what is
visible; the output settings decide how the picture is sampled and placed in the
document. The \type {eye} is the camera position, and \type {target} is the point
it looks at. The \type {up} direction rotates the view. In a perspective scene,
\type {fov} controls its width: a smaller value shows less and a larger value
shows more.

\type {crop} controls the relationship between the camera frame and the
objects. With \type {crop = true}, the result is fitted to the visible scene
bounds. With \type {crop = false}, the camera frame stays fixed and objects
outside it can be clipped. This makes it easier to compare settings such as
\type {fov} and \type {target}.

\index {raster resolution}%
At first, \type {width} and \type {height} are enough. They set the displayed
document dimensions. \type {resolution} converts one dimension to
raster pixels, and the other raster dimension is derived from the scene bounds.
Adding an object can therefore change the raster aspect ratio. If both document
dimensions are given, they fix the displayed rectangle.

\type {bytewidth} and \type {byteheight} choose the raster dimensions directly.
Give both byte dimensions and omit \type {width} and \type {height} when the
raster shape must stay fixed:

\starttyping[option=MP]
lmt_scene_start [
  bytewidth  = 800,
  byteheight = 500,
  projection = "perspective",
  eye        = (0,-8,3),
  target     = (0,0,0),
  fov        = 45,
  crop       = false,
] ;
% add the scene objects here
lmt_scene_stop ;
\stoptyping

This keeps the raster aspect ratio fixed. Use \type {width} or \type {height}
when physical size is the main concern; use an explicit byte-size pair when the
pixel dimensions must stay fixed.

\startsubsubject[title=Choosing an object]

Choose an object command from the way the geometry is described:

\starttabulate[|l|i1p(6cm)|i1p(4.5cm)|]
\NC \bf Object \NC \bf Use it when \NC \bf Command \NC \NR
\NC Graph surface \NC the height is given by \m {z=f(x,y)} \NC \goto {\type {lmt_scene_surface}} [surface-graphs] \NC \NR
\NC Parametric surface \NC the point \m {(x,y,z)} depends on \m {(u,v)} \NC \goto {\type {lmt_scene_parametric}} [parametric-surfaces] \NC \NR
\NC Implicit surface \NC an equation \m {F(x,y,z)=c} describes the shape \NC \goto {\type {lmt_scene_implicit}} [implicit-surfaces] \NC \NR
\NC Parametric curve \NC the point \m {(x,y,z)} depends on \m {t} \NC \goto {\type {lmt_scene_curve}} [parametric-curves] \NC \NR
\NC Plane \NC a finite rectangular patch is enough \NC \goto {\type {lmt_scene_plane}} [planes] \NC \NR
\NC External mesh \NC the triangles are read from an external file \NC \goto {\type {lmt_scene_external}} [external-meshes] \NC \NR
\NC Internal mesh \NC the geometry is generated by a \LUA\ mesher \NC \goto {\type {lmt_scene_internal}} [internal-meshes] \NC \NR
\stoptabulate

The first four are the usual starting points. Use a plane when a finite patch is
enough. Use an external or internal mesh when the geometry already exists in
another form or must be generated by \LUA.

\stopsubsubject

\startsubsubject[reference=defaults,,title=Defaults]

The following table collects the principal defaults in one place. A value marked
\type {unset} is derived from the scene bounds or from another dimension. Object
sampling defaults are included here because they determine geometric detail,
whereas \type {resolution} and \type {supersample} determine raster detail.

\starttabulate[|l|l|p(7cm)|]
\NC \bf Group \NC \bf Setting \NC \bf Default and meaning \NC \NR
\NC Scene \NC \type {width}, \type {height} \NC \type {unset}; the displayed size is derived unless a dimension is given. \NC \NR
\NC Scene \NC \type {bytemap} \NC integer slot \type {1}; use another slot for a separate scene output, or set it to \type {false} to disable bitmap output. \NC \NR
\NC Scene \NC \type {cache} \NC \type {false}; whole-scene rendered-bytemap caching is off. \NC \NR
\NC Camera \NC \type {projection} \NC \type {"perspective"}. \NC \NR
\NC Camera \NC \type {eye} \NC \type {(3,-4,2)}; \type {"auto"} fits a generated or imported scene. \NC \NR
\NC Camera \NC \type {target} \NC \type {(0,0,0)}; \type {"auto"} aims at the fitted scene center. \NC \NR
\NC Camera \NC \type {up}, \type {fov} \NC \type {(0,0,1)} and \type {45} degrees for perspective projection. \NC \NR
\NC Camera \NC \type {crop} \NC automatic scene-bound fitting when not explicitly set. \NC \NR
\NC Lighting \NC \type {light} \NC directional vector \type {(-1,-1,-1)}. \NC \NR
\NC Lighting \NC \type {intensity} \NC \type {1}. \NC \NR
\NC Lighting \NC \type {backgroundcolor} \NC white, \type {(1,1,1)}. \NC \NR
\NC Material \NC \type {diffuse} \NC \type {(.70,.75,.95)}. \NC \NR
\NC Material \NC \type {specular}, \type {shininess} \NC \type {(.50,.50,.50)} and \type {40}. \NC \NR
\NC Material \NC \type {ambient}, \type {opacity} \NC \type {.18} and \type {1}. \NC \NR
\NC Material \NC \type {twosided}, \type {ontop}, \type {dx}, \type {dy} \NC \type {false}, \type {false}, \type {0}, and \type {0}. \NC \NR
\NC Raster \NC \type {resolution} \NC \type {300} raster pixels per inch when byte dimensions are not explicit. \NC \NR
\NC Raster \NC \type {bytewidth}, \type {byteheight} \NC \type {unset}; derived from display dimensions and bounds. \NC \NR
\NC Raster \NC \type {maxfragment} \NC \type {0}; use the normal single-depth storage. \NC \NR
\NC Sampling \NC \type {supersample} \NC \type {1}; no temporary oversampling. \NC \NR
\NC Sampling \NC surface \type {nx}, \type {ny} \NC \type {32}, \type {32} intervals. \NC \NR
\NC Sampling \NC parametric \type {nu}, \type {nv} \NC \type {64}, \type {32} intervals. \NC \NR
\NC Sampling \NC curve \type {nt} \NC \type {64} samples. \NC \NR
\NC Sampling \NC implicit \type {nx}, \type {ny}, \type {nz} \NC required; there is no safe implicit-grid default. \NC \NR
\stoptabulate

\stopsubsubject

\startsubsubject[title=Scene keywords]

\startmixedcolumns

\keywordindex {bytemap}
\keyword {bytemap} selects the integer slot used for the rendered \METAPOST\
bitmap. The default is \type {1}; choose another slot when multiple scenes are
rendered in one picture and their bitmap outputs must remain distinct. Set
\type {bytemap = false} to disable bitmap output; this cannot be combined with
\type {cache = true}.

\keywordindex {name}
\keyword {name} gives a cached scene its name. Use a distinct name for each
scene definition that should have its own cache entry.

\keywordindex {cache}
\keyword {cache} when true enables the whole-scene rendered-bytemap cache. The
default is disabled; caching requires an enabled integer \type {bytemap}.

\keywordindex {width/height}
\keyword {width/height} sets the displayed width or height of the final result.
These are document dimensions, not necessarily the number of pixels used while
rendering. If only one is given, the other displayed dimension follows the
raster aspect ratio; giving both fixes the displayed rectangle.

\keywordindex {bytewidth}
\keyword {bytewidth} sets the raster width before scaling to the displayed size.
Use this together with \type {byteheight}, and without \type {width} or
\type {height}, when the pixel size and raster aspect ratio should be chosen
directly.

\keywordindex {byteheight}
\keyword {byteheight} sets the raster height before scaling to the displayed size.
Use this together with \type {bytewidth}, and without \type {width} or
\type {height}, when the pixel size and raster aspect ratio should be chosen
directly.

\keywordindex {supersample}
\keyword {supersample} sets the internal oversampling factor. It can make edges
smoother at the cost of a larger temporary raster. The maximum is \type {4}.

\keywordindex {maxfragment}
\keyword {maxfragment} limits the number of transparent depth fragments retained
per raster pixel. A larger value can preserve more overlapping translucent
surfaces, at the cost of more memory and compositing work. Values below one are
clamped internally; leave the default unless deep transparency is being lost.

\keywordindex {resolution}
\keyword {resolution} sets the conversion from a requested display size to raster
pixels when no explicit byte-size pair is given.

\keywordindex {backgroundcolor}
\keyword {backgroundcolor} sets the raster background color.

\keywordindex {projection}
\keyword {projection} selects the camera projection, normally \type {perspective}
or \type {orthographic}. Perspective gives nearby parts more visual weight;
orthographic keeps parallel directions parallel.

\keywordindex {eye}
\keyword {eye} sets the camera position, in the same coordinate system as the
objects in the scene. The special value \type {"auto"} derives a useful position
from the bounds of generated or imported geometry.

\keywordindex {target}
\keyword {target} sets the point toward which the camera looks. The special
value \type {"auto"} uses the center of the scene bounds.

\keywordindex {up}
\keyword {up} sets the upward camera direction, and therefore controls how the
view is rotated around the line from \type {eye} to \type {target}.

\keywordindex {fov}
\keyword {fov} sets the field of view for perspective projection. Smaller values
act more like a long lens; larger values show more of the surroundings.

\keywordindex {crop}
\keyword {crop} fits the rendered bitmap to the projected scene bounds when true.
When false, the camera frame is kept and the bitmap is not trimmed to the
objects.

\keywordindex {light}
\keyword {light} sets the direction of the directional light source. It is a
direction, not a point where a lamp is placed; \type {(0,0,-1)} means that light
travels down from above. The vector is normalized internally, so its length does
not change the result.

\keywordindex {intensity}
\keyword {intensity} sets the strength of the directional light source. The
default is \type {1}; lower values make the directional contribution weaker,
while the material's \type {ambient} value still provides a base level.

\keywordindex {stipple}
\keyword {stipple} enables the raster stipple postprocessor and passes it a
settings table. Its controls are described in the Stippling section.

\keywordindex {process}
\keyword {process} selects a named \LUA\ raster-postprocessing hook. The hook is
registered with \type {metapost.registermeshupper}; it runs on the rendered
raster before the supersampled result is resolved. See Raster postprocessing
hooks.

\stopmixedcolumns

\stopsubsubject

\startsubsubject[title=Scene settings examples]

Each example changes one scene setting. The geometry stays the same, so the
visible difference comes from the scene setup. The examples use
\type {crop = true} except the \type {fov} and \type {target} comparisons, which
keep the camera frame fixed.

\startbuffer[scene-reference]
lmt_scene_start [
  width = 6cm,
  crop  = true,
] ;

lmt_scene_surface [
  code = ".55*exp(-.07*(x*x+y*y))*(cos(2.2*x)+sin(2.0*y))",
  xmin = -3,
  xmax =  3,
  ymin = -3,
  ymax =  3,
] ;

lmt_scene_stop ;
\stopbuffer

\typebuffer[scene-reference][option=MP]

\startbuffer[scene-widthresolution]
lmt_scene_start [
  width = 6cm,
  crop  = true,
  resolution = 30,
] ;

lmt_scene_surface [
  code = ".55*exp(-.07*(x*x+y*y))*(cos(2.2*x)+sin(2.0*y))",
  xmin = -3,
  xmax =  3,
  ymin = -3,
  ymax =  3,
] ;

lmt_scene_stop ;
\stopbuffer

\typebuffer[scene-widthresolution][option=MP]

\startplacefloat
  [figure]
  [title={Showing the influence of the \type {resolution} keyword.}]
  \startcombination[nx=2,ny=1]
    {\startframed[TightBox]\processMPbuffer[scene-reference]\stopframed}{Default resolution 300dpi}
    {\startframed[TightBox]\processMPbuffer[scene-widthresolution]\stopframed}{\type {resolution = 30}}
  \stopcombination
\stopplacefloat

\startbuffer[scene-widthnofov]
lmt_scene_start [
  width = 6cm,
  crop  = false,
] ;

lmt_scene_surface [
  code = ".55*exp(-.07*(x*x+y*y))*(cos(2.2*x)+sin(2.0*y))",
  xmin = -3,
  xmax =  3,
  ymin = -3,
  ymax =  3,
] ;

lmt_scene_stop ;
\stopbuffer

\startbuffer[scene-widthfov]
lmt_scene_start [
  width = 6cm,
  fov   = 30,
  crop  = false,
] ;

lmt_scene_surface [
  code = ".55*exp(-.07*(x*x+y*y))*(cos(2.2*x)+sin(2.0*y))",
  xmin = -3,
  xmax =  3,
  ymin = -3,
  ymax =  3,
] ;

lmt_scene_stop ;
\stopbuffer

\startplacefloat
  [figure]
  [title={Setting the field of view. Its effect is easiest to see when
   \type {crop = false}.}]
  \startcombination[nx=2,ny=1]
    {\startframed[TightBox]\processMPbuffer[scene-widthnofov]\stopframed}{Default \type {fov = 45}}
    {\startframed[TightBox]\processMPbuffer[scene-widthfov]\stopframed}{\type {fov = 30}}
  \stopcombination
\stopplacefloat

\startbuffer[scene-widtheye]
lmt_scene_start [
  width = 6cm,
  eye   = (3,-4,4),
  crop  = true,
] ;

lmt_scene_surface [
  code = ".55*exp(-.07*(x*x+y*y))*(cos(2.2*x)+sin(2.0*y))",
  xmin = -3,
  xmax =  3,
  ymin = -3,
  ymax =  3,
] ;

lmt_scene_stop ;
\stopbuffer

\startplacefloat
  [figure]
  [title={Setting the position of the eye. Because the figure has a specified
   width, changing the eye can change its height substantially.}]
  \startcombination[nx=2,ny=1]
    {\startframed[TightBox]\processMPbuffer[scene-reference]\stopframed}{Default \type {eye = (3,-4,2)}}
    {\startframed[TightBox]\processMPbuffer[scene-widtheye]\stopframed}{\type {eye = (3,-4,4)}}
  \stopcombination
\stopplacefloat

\startbuffer[scene-widthtarget]
lmt_scene_start [
  width  = 6cm,
  target = (1,0,0),
  crop   = false,
] ;

lmt_scene_surface [
  code = ".55*exp(-.07*(x*x+y*y))*(cos(2.2*x)+sin(2.0*y))",
  xmin = -3,
  xmax =  3,
  ymin = -3,
  ymax =  3,
] ;

lmt_scene_stop ;
\stopbuffer

\startplacefloat
  [figure]
  [title={Setting the target point. Its effect is easiest to see when
   \type {crop = false}.}]
  \startcombination[nx=2,ny=1]
    {\startframed[TightBox]\processMPbuffer[scene-widthnofov]\stopframed}{Default \type {target = (0,0,0)}}
    {\startframed[TightBox]\processMPbuffer[scene-widthtarget]\stopframed}{\type {target = (1,0,0)}}
  \stopcombination
\stopplacefloat

\startbuffer[scene-up-xaxis]
lmt_scene_start [
  width = 6cm,
  up    = (1,0,0),
  crop  = true,
] ;

lmt_scene_surface [
  code = ".55*exp(-.07*(x*x+y*y))*(cos(2.2*x)+sin(2.0*y))",
  xmin = -3,
  xmax =  3,
  ymin = -3,
  ymax =  3,
] ;

lmt_scene_stop ;
\stopbuffer

\startplacefloat
  [figure]
  [title={Setting the \type {up} direction rolls the camera frame around the
   eye--target axis.}]
  \startcombination[nx=2,ny=1]
    {\startframed[TightBox]\processMPbuffer[scene-reference]\stopframed}{Default \type {up = (0,0,1)}}
    {\startframed[TightBox]\processMPbuffer[scene-up-xaxis]\stopframed}{\type {up = (1,0,0)}}
  \stopcombination
\stopplacefloat

\startbuffer[scene-widthlowres]
lmt_scene_start [
  width      = 6cm,
  resolution = 30,
  crop       = true,
] ;

lmt_scene_surface [
  code = ".55*exp(-.07*(x*x+y*y))*(cos(2.2*x)+sin(2.0*y))",
  xmin = -3,
  xmax =  3,
  ymin = -3,
  ymax =  3,
] ;

lmt_scene_stop ;
\stopbuffer

\startbuffer[scene-widthsupersample]
lmt_scene_start [
  width       = 6cm,
  resolution  = 30,
  crop        = true,
  supersample = 4,
] ;

lmt_scene_surface [
  code = ".55*exp(-.07*(x*x+y*y))*(cos(2.2*x)+sin(2.0*y))",
  xmin = -3,
  xmax =  3,
  ymin = -3,
  ymax =  3,
] ;

lmt_scene_stop ;
\stopbuffer

\startplacefloat
  [figure]
  [title={Supersampling produces antialiasing. In this example, we also use a low
   resolution to make the difference clear.}]
  \startcombination[nx=2,ny=1]
    {\startframed[TightBox]\processMPbuffer[scene-widthlowres]\stopframed}{Default}
    {\startframed[TightBox]\processMPbuffer[scene-widthsupersample]\stopframed}{\type {supersample = 4}}
  \stopcombination
\stopplacefloat

\startbuffer[scene-widthlight]
lmt_scene_start [
  width = 6cm,
  crop  = true,
  light = (0,0,-1),
] ;

lmt_scene_surface [
  code = ".55*exp(-.07*(x*x+y*y))*(cos(2.2*x)+sin(2.0*y))",
  xmin = -3,
  xmax =  3,
  ymin = -3,
  ymax =  3,
] ;

lmt_scene_stop ;
\stopbuffer

\startplacefloat
  [figure]
  [title={Changing the direction of the light. The vector is a direction, not a
   position; \type {(0,0,-1)} points down from above.}]
  \startcombination[nx=2,ny=1]
    {\startframed[TightBox]\processMPbuffer[scene-reference]\stopframed}{Default \type {light = (-1,-1,-1)}}
    {\startframed[TightBox]\processMPbuffer[scene-widthlight]\stopframed}{\type {light = (0,0,-1)}}
  \stopcombination
\stopplacefloat

\startbuffer[scene-widthintensity]
lmt_scene_start [
  width     = 6cm,
  crop      = true,
  intensity = 0.5,
] ;

lmt_scene_surface [
  code = ".55*exp(-.07*(x*x+y*y))*(cos(2.2*x)+sin(2.0*y))",
  xmin = -3,
  xmax =  3,
  ymin = -3,
  ymax =  3,
] ;

lmt_scene_stop ;
\stopbuffer

\startplacefloat
  [figure]
  [title={The \type {intensity} key sets the strength of the light.}]
  \startcombination[nx=2,ny=1]
    {\startframed[TightBox]\processMPbuffer[scene-reference]\stopframed}{Default \type {intensity = 1}}
    {\startframed[TightBox]\processMPbuffer[scene-widthintensity]\stopframed}{\type {intensity = 0.5}}
  \stopcombination
\stopplacefloat

\startbuffer[scene-widthorthographic]
lmt_scene_start [
  width      = 6cm,
  crop       = true,
  projection = "orthographic",
] ;

lmt_scene_surface [
  code = ".55*exp(-.07*(x*x+y*y))*(cos(2.2*x)+sin(2.0*y))",
  xmin = -3,
  xmax =  3,
  ymin = -3,
  ymax =  3,
] ;

lmt_scene_stop ;
\stopbuffer

\startplacefloat
  [figure]
  [title={The \type {projection} keyword selects the projection type. The default
   is \type {projection = "perspective"}.}]
  \startcombination[nx=2,ny=1]
    {\startframed[TightBox]\processMPbuffer[scene-reference]\stopframed}{\type {projection = "perspective"}}
    {\startframed[TightBox]\processMPbuffer[scene-widthorthographic]\stopframed}{\type {projection = "orthographic"}}
  \stopcombination
\stopplacefloat

\startbuffer[scene-widthbackground]
lmt_scene_start [
  width           = 6cm,
  crop            = true,
  backgroundcolor = (0.7,0.1,0.1),
] ;

lmt_scene_surface [
  code = ".55*exp(-.07*(x*x+y*y))*(cos(2.2*x)+sin(2.0*y))",
  xmin = -3,
  xmax =  3,
  ymin = -3,
  ymax =  3,
] ;

lmt_scene_stop ;
\stopbuffer

\startplacefloat
  [figure]
  [title={Changing the background color with \type {backgroundcolor}.}]
  \startcombination[nx=2,ny=1]
    {\startframed[TightBox]\processMPbuffer[scene-reference]\stopframed}{Default, not set}
    {\startframed[TightBox]\processMPbuffer[scene-widthbackground]\stopframed}{\type {backgroundcolor = (0.7,0.1,0.1)}}
  \stopcombination
\stopplacefloat

\stopsubsubject

\stopsection

\startsection[reference=materials,title=Materials]

A material describes how a rendered object looks when the scene light reaches
it. It is separate from the geometry itself: a surface, mesh, curve, or generated
object can keep the same shape while getting a different color, highlight,
ambient level, or transparency. When no material is selected, the scene uses the
default material. Named materials are useful when several objects should share
the same look.

We define a material with

\starttyping[option=MP]
lmt_scene_material [ ... ]
\stoptyping

A material might look like this:

\starttyping[option=MP]
lmt_scene_material [
  name      = "softgreen",
  diffuse   = (.1,.8,.1),
  specular  = (.25,.25,.5),
  shininess = 10,
  ambient   = .28,
  opacity   = 1,
] ;
\stoptyping

The renderer uses Lambertian diffuse shading with a Blinn--Phong specular term.
The diffuse part follows the cosine rule for a matte surface. The optional
highlight uses the halfway direction between the light and the eye. For the
historical background, see the original
\goto {Phong paper} [url(https://doi.org/10.1145/360825.360839)] and
\goto {Blinn's refinement} [url(https://doi.org/10.1145/563858.563896)]. This is
an illustration model, not a physically based renderer: the light is
directional, there is no distance falloff or cast shadow, and final colors are
clipped to the displayable range.

The public \type {light} vector gives the direction in which light travels, not
the position of a lamp. Thus \type {(0,0,-1)} means that light comes from above
and travels downwards. Its length does not matter. A surface is brightest when
its normal points toward the light and darkest when it points away. The
\type {ambient} value provides a floor, so the dark side need not be black. In
rough terms, diffuse brightness is the ambient value plus an angle-dependent
contribution, multiplied by \type {intensity} and limited to \type {0}--\type {1}.
With \type {twosided = true}, the back side is lit as well.

\index {opacity}%
The \RGB\ \type {diffuse} triplet is the object's color. The renderer modulates it
by the diffuse brightness and then adds the highlight from \type {specular}.
Set \type {specular = (0,0,0)} for a matte object. Otherwise, \type {shininess}
controls the size of the highlight: smaller values make it broad, while larger
values make it small and sharp. The light \type {intensity} affects both the
diffuse contribution and the highlight.

Start by choosing \type {diffuse}, keep \type {ambient} near its default, and
adjust \type {light} if the shape is hard to read. Add a modest \type {specular}
value for a glossy surface and increase \type {shininess} if the highlight is
too broad.

\startbuffer[material-ambient]
lmt_scene_start [
  width = 6cm,
  crop  = true,
] ;

lmt_scene_material [
  name      = "ambienthill",
  diffuse   = (.70,.75,.95),
  specular  = (.50,.50,.50),
  shininess = 40,
  ambient   = .5,
] ;

lmt_scene_surface [
  code     = ".55*exp(-.07*(x*x+y*y))*(cos(2.2*x)+sin(2.0*y))",
  xmin     = -3,
  xmax     =  3,
  ymin     = -3,
  ymax     =  3,
  material = "ambienthill",
] ;

lmt_scene_stop ;
\stopbuffer

\startplacefloat
  [figure]
  [title={The material's \type {ambient} value lifts the dark side without changing
   the direction of the light.}]
  \startcombination[nx=2,ny=1]
    {\startframed[TightBox]\processMPbuffer[scene-reference]\stopframed}{Default \type {ambient = .18}}
    {\startframed[TightBox]\processMPbuffer[material-ambient]\stopframed}{\type {ambient = .5}}
  \stopcombination
\stopplacefloat

\startbuffer[material-matte]
lmt_scene_start [
  width = 6cm,
  crop  = true,
] ;

lmt_scene_material [
  name      = "mattehill",
  diffuse   = (.70,.75,.95),
  specular  = (0,0,0),
  shininess = 40,
  ambient   = .18,
] ;

lmt_scene_surface [
  code     = ".55*exp(-.07*(x*x+y*y))*(cos(2.2*x)+sin(2.0*y))",
  xmin     = -3,
  xmax     =  3,
  ymin     = -3,
  ymax     =  3,
  material = "mattehill",
] ;

lmt_scene_stop ;
\stopbuffer

\startbuffer[material-glossy]
lmt_scene_start [
  width = 6cm,
  crop  = true,
] ;

lmt_scene_material [
  name      = "glossyhill",
  diffuse   = (.70,.75,.95),
  specular  = (.50,.50,.50),
  shininess = 8,
  ambient   = .18,
] ;

lmt_scene_surface [
  code     = ".55*exp(-.07*(x*x+y*y))*(cos(2.2*x)+sin(2.0*y))",
  xmin     = -3,
  xmax     =  3,
  ymin     = -3,
  ymax     =  3,
  material = "glossyhill",
] ;

lmt_scene_stop ;
\stopbuffer

\startplacefloat
  [figure]
  [title={A \type {specular} color adds a highlight; \type {shininess} controls its
   size.}]
  \startcombination[nx=2,ny=1]
    {\startframed[TightBox]\processMPbuffer[material-matte]\stopframed}{matte}
    {\startframed[TightBox]\processMPbuffer[material-glossy]\stopframed}{glossy, \type {shininess = 8}}
  \stopcombination
\stopplacefloat

\index {line widths}%
An object selects this material with \type {material = "softgreen"}. The remaining
settings control reuse, lighting response, sidedness, transparency, and the
raster size of point-like or line-like meshes.

If no material is selected, the scene uses the opaque blue-gray default material
listed in the \goto {Defaults} [defaults] table. Defining a named material is
therefore only necessary when the default appearance is not what you want or
when several objects should share another appearance.

The material keywords are:

\startmixedcolumns
\keywordindex {name}
\keyword {name} gives the material a name so that surfaces and other objects can
refer to it with their \type {material} keyword.

\keywordindex {diffuse}
\keyword {diffuse} sets the main color that is affected by the directional light.
Choose this color first.

\keywordindex {specular}
\keyword {specular} sets the color of the specular highlight, the small bright
part that can appear where the light reflects toward the eye.

\keywordindex {shininess}
\keyword {shininess} sets the sharpness of the specular highlight; larger values
give a tighter highlight.

\keywordindex {ambient}
\keyword {ambient} sets the base light level that remains visible even where the
directional light contributes little or nothing. Higher values make shaded parts
less dark.

\keywordindex {twosided}
\keyword {twosided} also illuminates the back side of an open surface. With
\type {false}, a back-facing side receives only ambient light; with \type {true},
its normal is flipped for the lighting calculation. This helps with thin sheets
or open surfaces viewed from either side.

\keywordindex {opacity}
\keyword {opacity} sets the transparency of the material; \type {1} is fully
opaque and \type {0} is fully transparent. Values in between are rendered in a
separate transparent pass. A value around \type {.2} or \type {.3} is a useful
starting point for a cutting plane.

\keywordindex {ontop}
\keyword {ontop} renders the material in the final on-top pass. It can be useful
for helper geometry that should remain visible over the regular scene.

\keywordindex {dx}
\keyword {dx} sets the horizontal thickness of curve, intersection, and other
line-like or point-like material in raster pixels. For ordinary triangle
surfaces this is normally left at the default.

\keywordindex {dy}
\keyword {dy} sets the vertical thickness of curve, intersection, and other
line-like or point-like material in raster pixels. For ordinary triangle
surfaces this is normally left at the default.
\stopmixedcolumns

Opacity is easiest to judge against geometry behind the material. The two scenes
below keep the hill, plane, camera, and lighting fixed and change only the plane's
\type {opacity}.

\startbuffer[material-opacity-opaque]
lmt_scene_start [
  width  = 6cm,
  crop   = true,
  eye    = (3,-4,2.6),
  target = (0,0,.1),
  light  = (-1,-1,-2),
] ;

lmt_scene_material [
  name    = "opacityhill",
  diffuse = (.25,.55,.85),
  ambient = .25,
] ;

lmt_scene_material [
  name     = "opaqueplane",
  diffuse  = (.9,.25,.12),
  opacity  = 1,
  twosided = true,
] ;

lmt_scene_surface [
  code     = ".6*exp(-.6*(x*x+y*y)) - .2",
  xmin     = -1.8,
  xmax     =  1.8,
  ymin     = -1.8,
  ymax     =  1.8,
  material = "opacityhill",
] ;

lmt_scene_plane [
  center   = (0,0,.15),
  normal   = (0,0,1),
  usize    = 3,
  vsize    = 3,
  material = "opaqueplane",
] ;

lmt_scene_stop ;
\stopbuffer

\startbuffer[material-opacity-transparent]
lmt_scene_start [
  width  = 6cm,
  crop   = true,
  eye    = (3,-4,2.6),
  target = (0,0,.1),
  light  = (-1,-1,-2),
] ;

lmt_scene_material [
  name    = "opacityhill",
  diffuse = (.25,.55,.85),
  ambient = .25,
] ;

lmt_scene_material [
  name     = "transparentplane",
  diffuse  = (.9,.25,.12),
  opacity  = .25,
  twosided = true,
] ;

lmt_scene_surface [
  code     = ".6*exp(-.6*(x*x+y*y)) - .2",
  xmin     = -1.8,
  xmax     =  1.8,
  ymin     = -1.8,
  ymax     =  1.8,
  material = "opacityhill",
] ;

lmt_scene_plane [
  center   = (0,0,.15),
  normal   = (0,0,1),
  usize    = 3,
  vsize    = 3,
  material = "transparentplane",
] ;

lmt_scene_stop ;
\stopbuffer

\startplacefloat
  [figure]
  [title={The same cutting plane at two opacity values.}]
  \startcombination[nx=2,ny=1]
    {\processMPbuffer[material-opacity-opaque]}{\type {opacity = 1}}
    {\processMPbuffer[material-opacity-transparent]}{\type {opacity = .25}}
  \stopcombination
\stopplacefloat

\startsubsubject[title=Two-sided materials]

The effect of \type {twosided} is easiest to see on an open surface viewed from
the back. The scene below shows two horizontal patches from underneath.
The light is below them, so the patch with \type {twosided = false} receives no
directional light, while the other patch is lit after its normal is flipped.

\startbuffer[material-twosided]
lmt_scene_start [
  width  = 10cm,
  crop   = true,
  eye    = (4,-6,-4),
  target = (0,0,0),
  up     = (0,0,1),
  light  = (0,0,1),
] ;

lmt_scene_material [
  name     = "backface",
  diffuse  = (.15,.45,.85),
  specular = (0,0,0),
  ambient  = 0,
  twosided = false,
] ;

lmt_scene_material [
  name     = "twoside",
  diffuse  = (.15,.45,.85),
  specular = (0,0,0),
  ambient  = 0,
  twosided = true,
] ;

lmt_scene_plane [
  center   = (-1.5,0,0),
  normal   = (0,0,1),
  usize    = 2.4,
  vsize    = 2.4,
  material = "backface",
] ;

lmt_scene_plane [
  center   = (1.5,0,0),
  normal   = (0,0,1),
  usize    = 2.4,
  vsize    = 2.4,
  material = "twoside",
] ;

lmt_scene_stop ;
\stopbuffer

\typebuffer[material-twosided][option=MP]

\startplacefloat
  [figure]
  [title={Viewing two horizontal patches from below: \type {twosided = false}
   receives no directional light, while \type {twosided = true} lights the back
   side.}]
  \startframed[TightBox]
    \processMPbuffer[material-twosided]
  \stopframed
\stopplacefloat

The two patches have the same geometry and diffuse color. The difference comes
only from the back-face lighting rule; setting \type {twosided} does not make an
opaque surface transparent.

\stopsubsubject

\stopsection

\startsection[reference=surface-graphs,title=Surface graphs]

A surface graph is the 3D version of a familiar function plot. For each point
\m {(x,y)} in a chosen rectangle, the expression gives a height \m {z}. The
renderer samples that rectangle, builds a triangle mesh from the sample points,
and then shades the mesh like the other objects in the scene.

Add a surface graph \m {z = f(x,y)} with

\starttyping[option=MP]
lmt_scene_surface [ ... ]
\stoptyping

A graph surface might look like this:

\starttyping[option=MP]
lmt_scene_surface [
  code = "sin(x)*cos(y)",
  xmin = -pi,
  xmax =  pi,
  ymin = -pi,
  ymax =  pi,
  nx   = 32,
  ny   = 32,
] ;
\stoptyping

The \type {code} value gives the height. The four bounds select the part of the
\m {(x,y)} plane to sample, and \type {nx} and \type {ny} control the mesh density.
These are the only object settings needed for a simple graph. The other settings
handle naming, lighting normals, and material selection.

\index {normals}%
A normal is a direction perpendicular to the surface. It tells the lighting
model which way each part of the mesh faces. Without derivative expressions, the
renderer estimates normals from the sampled mesh. The \type {normal} setting
chooses how to use them: \type {smooth} gives a continuous appearance,
\type {flat} shows the triangle faces, and \type {up} uses the same upward
direction everywhere.

The graph keywords are:

\startmixedcolumns
\keywordindex {id}
\keyword {id} gives the surface a name so that it can be referred to elsewhere,
for example by an intersection.

\keywordindex {code}
\keyword {code} gives the function \m {z=f(x,y)} as an expression in \type {x} and
\type {y}. It is evaluated at each sampled point in the domain.

\keywordindex {dfdx}
\keyword {dfdx} gives the slope in the \type {x} direction. It is optional; supply
it together with \type {dfdy} when analytic slope information is available to
improve the lighting normals.

\keywordindex {dfdy}
\keyword {dfdy} gives the slope in the \type {y} direction. Together with
\type {dfdx}, it can be used when calculating more accurate surface normals.

\keywordindex {normal}
\keyword {normal} selects how normals are made for lighting. \type {smooth}
interpolates normals across the mesh, \type {flat} exposes the triangle faces, and
\type {up} uses \type {(0,0,1)} everywhere.

\keywordindex {xmin}
\keyword {xmin} sets the lower bound of the sampled rectangle in the \type {x}
direction.

\keywordindex {xmax}
\keyword {xmax} sets the upper bound of the sampled rectangle in the \type {x}
direction.

\keywordindex {ymin}
\keyword {ymin} sets the lower bound of the sampled rectangle in the \type {y}
direction.

\keywordindex {ymax}
\keyword {ymax} sets the upper bound of the sampled rectangle in the \type {y}
direction.

\keywordindex {nx}
\keyword {nx} sets the number of mesh intervals in the \type {x} direction. Higher
values give more detail and more triangles.

\keywordindex {ny}
\keyword {ny} sets the number of mesh intervals in the \type {y} direction. Higher
values give more detail and more triangles.

\keywordindex {material}
\keyword {material} selects the named material used when rendering the surface.

\stopmixedcolumns

The examples below keep the function fixed and change one choice at a time:
mesh resolution, normal appearance, slope information, material, and cropping.
That makes the effect of each setting easier to see.

\startbuffer[surface]
lmt_scene_start [
  width      = 6cm,
  crop      = true,
] ;

lmt_scene_surface [
  code = ".55*exp(-.07*(x*x+y*y))*(cos(2.2*x)+sin(2.0*y))",
  xmin = -3,
  xmax =  3,
  ymin = -3,
  ymax =  3,
] ;

lmt_scene_stop ;
\stopbuffer

\startbuffer[surface-a]
lmt_scene_start [
  width      = 6cm,
  crop      = true,
] ;

lmt_scene_surface [
  code = ".55*exp(-.07*(x*x+y*y))*(cos(2.2*x)+sin(2.0*y))",
  xmin = -3,
  xmax =  3,
  ymin = -3,
  ymax =  3,
  nx   = 10,
  ny   = 10,
] ;

lmt_scene_stop ;
\stopbuffer

\typebuffer[surface][option=MP]

\startplacefloat
  [figure]
  [title={Changing the graph mesh interval counts.}]
  \startcombination[nx=2,ny=1]
    {\processMPbuffer[surface]}{default \type {nx = 32, ny = 32}}
    {\processMPbuffer[surface-a]}{\type {nx = 10, ny = 10}}
  \stopcombination
\stopplacefloat

\startbuffer[surface-normal-smooth]
lmt_scene_start [
  width      = 4cm,
  crop      = true,
] ;

lmt_scene_surface [
  code = ".55*exp(-.07*(x*x+y*y))*(cos(2.2*x)+sin(2.0*y))",
  xmin = -3,
  xmax =  3,
  ymin = -3,
  ymax =  3,
  normal = "smooth",
] ;

lmt_scene_stop ;
\stopbuffer

\startbuffer[surface-normal-flat]
lmt_scene_start [
  width      = 4cm,
  crop      = true,
] ;

lmt_scene_surface [
  code = ".55*exp(-.07*(x*x+y*y))*(cos(2.2*x)+sin(2.0*y))",
  xmin = -3,
  xmax =  3,
  ymin = -3,
  ymax =  3,
  normal = "flat",
] ;

lmt_scene_stop ;
\stopbuffer

\startbuffer[surface-normal-up]
lmt_scene_start [
  width      = 4cm,
  crop      = true,
] ;

lmt_scene_surface [
  code = ".55*exp(-.07*(x*x+y*y))*(cos(2.2*x)+sin(2.0*y))",
  xmin = -3,
  xmax =  3,
  ymin = -3,
  ymax =  3,
  normal = "up",
] ;

lmt_scene_stop ;
\stopbuffer

\startplacefloat
  [figure]
  [title={Different normal modes. \type {normal = "smooth"} is the default.}]
  \startcombination[nx=3,ny=1]
    {\processMPbuffer[surface-normal-smooth]}{\type {normal = "smooth"}}
    {\processMPbuffer[surface-normal-flat]}  {\type {normal = "flat"}}
    {\processMPbuffer[surface-normal-up]}    {\type {normal = "up"}}
  \stopcombination
\stopplacefloat

\startbuffer[surface-noderivatives]
lmt_scene_start [
  width      = 6cm,
  crop      = true,
] ;

lmt_scene_surface [
  code = ".55*exp(-.07*(x*x+y*y))*(cos(2.2*x)+sin(2.0*y))",
  % dfdx = ".55*exp(-.07*(x*x+y*y))*(-.14*x*(cos(2.2*x)+sin(2.0*y))-2.2*sin(2.2*x))",
  % dfdy = ".55*exp(-.07*(x*x+y*y))*(-.14*y*(cos(2.2*x)+sin(2.0*y))+2.0*cos(2.0*y))",
  xmin = -3,
  xmax =  3,
  ymin = -3,
  ymax =  3,
  nx   = 10,
  ny   = 10,
] ;

lmt_scene_stop ;
\stopbuffer

\startbuffer[surface-derivatives]
lmt_scene_start [
  width      = 6cm,
  crop      = true,
] ;

lmt_scene_surface [
  code = ".55*exp(-.07*(x*x+y*y))*(cos(2.2*x)+sin(2.0*y))",
  dfdx = ".55*exp(-.07*(x*x+y*y))*(-.14*x*(cos(2.2*x)+sin(2.0*y))-2.2*sin(2.2*x))",
  dfdy = ".55*exp(-.07*(x*x+y*y))*(-.14*y*(cos(2.2*x)+sin(2.0*y))+2.0*cos(2.0*y))",
  xmin = -3,
  xmax =  3,
  ymin = -3,
  ymax =  3,
  nx   = 10,
  ny   = 10,
] ;

lmt_scene_stop ;
\stopbuffer

\startplacefloat
  [figure]
  [title={Comparing a surface with and without slope expressions.}]
  \startcombination[nx=2,ny=1]
    {\processMPbuffer[surface-noderivatives]}{\type {nx = 10, ny = 10}}
    {\processMPbuffer[surface-derivatives]}  {also \type {dfdx = ..., dfdy = ...}}
  \stopcombination
\stopplacefloat

\startbuffer[surface-material]
lmt_scene_start [
  width      = 6cm,
  crop      = true,
] ;

lmt_scene_material [
    name      = "mymaterial",
    diffuse   = (.1,.8,.1),
    specular  = (.25,.25,.5),
    shininess = 10,
    ambient   = .28,
    opacity   = 1,
] ;

lmt_scene_surface [
  code = ".55*exp(-.07*(x*x+y*y))*(cos(2.2*x)+sin(2.0*y))",
  xmin = -3,
  xmax =  3,
  ymin = -3,
  ymax =  3,
  material = "mymaterial",
] ;

lmt_scene_stop ;
\stopbuffer

Change a surface's appearance by defining a material and selecting it with the
\type {material} keyword.

\starttyping[option=MP]
lmt_scene_material [
    name      = "mymaterial",
    diffuse   = (.1,.8,.1),
    specular  = (.25,.25,.5),
    shininess = 10,
    ambient   = .28,
    opacity   = 1,
] ;
\stoptyping

\startplacefloat
  [figure]
  [title={Changing a surface's appearance with a material.}]
  \startcombination[nx=2,ny=1]
    {\processMPbuffer[surface]}{default}
    {\processMPbuffer[surface-material]}  {\type {material = "mymaterial"}}
  \stopcombination
\stopplacefloat

\startbuffer[surface-cropfalse]
lmt_scene_start [
  width      = 6cm,
  crop      = false,
  fov        = 27,
  eye        = (6,-8,4),
] ;

lmt_scene_surface [
  code = ".55*exp(-.07*(x*x+y*y))*(cos(2.2*x)+sin(2.0*y))",
  xmin = -3,
  xmax =  3,
  ymin = -3,
  ymax =  3,
] ;

lmt_scene_stop ;
\stopbuffer

\startplacefloat
  [figure]
  [title={Cropping the rendered result.}]
  \startcombination[nx=2,ny=1]
    {\startframed[TightBox]\processMPbuffer[surface]\stopframed}          {default}
    {\startframed[TightBox]\processMPbuffer[surface-cropfalse]\stopframed}{\type {crop = false}}
  \stopcombination
\stopplacefloat
\stopsection

\startsection[reference=parametric-surfaces,title=Parametric surfaces]

A parametric surface is built from two parameters rather than from a height
function. For each sampled pair \m {(u,v)}, three expressions give a point
\m {(x,y,z)} in space. This covers shapes that are not naturally described as
\m {z=f(x,y)}, such as tori, spheres, tubes, and twisted patches.

Add a parametric surface with

\starttyping[option=MP]
lmt_scene_parametric [ ... ]
\stoptyping

A parametric surface might look like this:

\starttyping[option=MP]
lmt_scene_parametric [
  x    = "(1+.35*cos(v))*cos(u)",
  y    = "(1+.35*cos(v))*sin(u)",
  z    = ".35*sin(v)",
  umin = 0,
  umax = 2*pi,
  vmin = 0,
  vmax = 2*pi,
  nu   = 64,
  nv   = 32,
] ;
\stoptyping

\index {normals}%
The coordinate expressions describe the shape. The \type{u} and \type{v} bounds
choose the part of the parameter plane to sample, and \type{nu} and \type{nv}
control the mesh resolution. The parameters \type {u} and \type {v} are just
variables; they need not be spatial coordinates. The other settings handle
naming, lighting normals, and material selection.

The parametric-surface keywords are:

\startmixedcolumns

\keywordindex {id}
\keyword {id} gives the surface a name so that it can be referred to elsewhere,
for example by an intersection.

\keywordindex {x}
\keyword {x} gives the \type{x(u,v)} coordinate as an expression in \type{u} and
\type{v}. It is evaluated at each sampled parameter pair.

\keywordindex {y}
\keyword {y} gives the \type{y(u,v)} coordinate as an expression in \type{u} and
\type{v}. It is evaluated at each sampled parameter pair.

\keywordindex {z}
\keyword {z} gives the \type{z(u,v)} coordinate as an expression in \type{u} and
\type{v}. It is evaluated at each sampled parameter pair.

\keywordindex {xu/yu/zu}
\keyword {xu/yu/zu} describe the tangent direction when \type{u} changes. They
are optional; together with the other three tangent components they provide
analytic normals.

\keywordindex {xv/yv/zv}
\keyword {xv/yv/zv} describe the tangent direction when \type{v} changes. Give
all six derivative expressions together when analytic normals are desired. If
they are omitted, normals are estimated from the sampled triangles.

\keywordindex {normal}
\keyword {normal} selects how normals are made for lighting. \type{smooth}
interpolates normals across the mesh, \type{flat} exposes the triangle faces, and
\type{up} uses \type{(0,0,1)} everywhere.

\keywordindex {normalstep}
\keyword {normalstep} sets the finite-difference step used when analytic tangent
expressions are not supplied. A positive value can stabilize numerical normals
for very small or very large parameter domains; otherwise an automatic
domain-scaled step is used.

\keywordindex {umin}
\keyword {umin} sets the lower bound of the sampled rectangle in the \type{u}
direction.

\keywordindex {umax}
\keyword {umax} sets the upper bound of the sampled rectangle in the \type{u}
direction.

\keywordindex {vmin}
\keyword {vmin} sets the lower bound of the sampled rectangle in the \type{v}
direction.

\keywordindex {vmax}
\keyword {vmax} sets the upper bound of the sampled rectangle in the \type{v}
direction.

\keywordindex {nu}
\keyword {nu} sets the number of mesh intervals in the \type{u} direction. Higher
values give more detail and more triangles.

\keywordindex {nv}
\keyword {nv} sets the number of mesh intervals in the \type{v} direction. Higher
values give more detail and more triangles.

\keywordindex {material}
\keyword {material} selects the named material used when rendering the surface.

\stopmixedcolumns

The examples below keep the torus formula fixed and change the mesh resolution,
derivative information, or normal mode.

\startbuffer[parametric]
lmt_scene_start [
  width      = 5cm,
  crop      = true,
] ;

lmt_scene_parametric [
  x = "(1+.35*cos(v))*cos(u)",
  y = "(1+.35*cos(v))*sin(u)",
  z = ".35*sin(v)",
] ;

lmt_scene_stop ;
\stopbuffer

\startbuffer[parametric-lowmesh]
lmt_scene_start [
  width      = 5cm,
  crop      = true,
] ;

lmt_scene_parametric [
  x  = "(1+.35*cos(v))*cos(u)",
  y  = "(1+.35*cos(v))*sin(u)",
  z  = ".35*sin(v)",
  nu = 16,
  nv = 8,
] ;

lmt_scene_stop ;
\stopbuffer

\typebuffer[parametric][option=MP]

\startplacefloat
  [figure]
  [title={Changing the parametric mesh resolution.}]
  \startcombination[nx=2,ny=1]
    {\processMPbuffer[parametric]}         {default}
    {\processMPbuffer[parametric-lowmesh]} {\type {nu = 16, nv = 8}}
  \stopcombination
\stopplacefloat

\startbuffer[parametric-derivatives]
lmt_scene_start [
  width      = 5cm,
  crop      = true,
] ;

lmt_scene_parametric [
  x  = "(1+.35*cos(v))*cos(u)",
  y  = "(1+.35*cos(v))*sin(u)",
  z  = ".35*sin(v)",
  xu = "-(1+.35*cos(v))*sin(u)",
  yu =  "(1+.35*cos(v))*cos(u)",
  zu = "0",
  xv = "-.35*sin(v)*cos(u)",
  yv = "-.35*sin(v)*sin(u)",
  zv =  ".35*cos(v)",
%   nu = 16,
%   nv = 8,
] ;

lmt_scene_stop ;
\stopbuffer

\startplacefloat
  [figure]
  [title={Giving derivatives for the parametric normals.}]
  \startcombination[nx=2,ny=1]
    {\processMPbuffer[parametric]}            {no derivatives}
    {\processMPbuffer[parametric-derivatives]}{with derivatives}
  \stopcombination
\stopplacefloat

\startbuffer[parametric-flat]
lmt_scene_start [
  width      = 5cm,
  crop      = true,
] ;

lmt_scene_parametric [
  x      = "(1+.35*cos(v))*cos(u)",
  y      = "(1+.35*cos(v))*sin(u)",
  z      = ".35*sin(v)",
  normal = "flat",
] ;

lmt_scene_stop ;
\stopbuffer

\startbuffer[parametric-up]
lmt_scene_start [
  width      = 5cm,
  crop      = true,
] ;

lmt_scene_parametric [
  x      = "(1+.35*cos(v))*cos(u)",
  y      = "(1+.35*cos(v))*sin(u)",
  z      = ".35*sin(v)",
  normal = "up",
] ;

lmt_scene_stop ;
\stopbuffer

\startplacefloat
  [figure]
  [title={Changing the normal mode of a parametric surface.}]
  \startcombination[nx=3,ny=1]
    {\processMPbuffer[parametric]}      {default}
    {\processMPbuffer[parametric-flat]} {\type {normal = "flat"}}
    {\processMPbuffer[parametric-up]}   {\type {normal = "up"}}
  \stopcombination
\stopplacefloat

\stopsection

\startsection[reference=implicit-surfaces,title=Implicit surfaces]

An implicit surface is described by an equation instead of explicit coordinates.
The usual form is \m {F(x,y,z)=c}: all points where the function has the chosen
value belong to the surface. This suits shapes that are more natural to describe
as a condition in space, for example a sphere with
\m {x^2+y^2+z^2=1}.

\index {normals}%
The renderer samples a finite 3D box and builds a triangle mesh where the
function crosses the chosen iso value. This is a local approximation, not a
global solution of the equation. The box and sampling resolution matter: parts
outside the box are not seen, and a feature can be missed when the grid is too
coarse to reveal it.

Add an implicit surface with

\starttyping[option=MP]
lmt_scene_implicit [ ... ]
\stoptyping

An implicit surface might look like this:

\starttyping[option=MP]
lmt_scene_implicit [
  code = "x*x + y*y + z*z - 1",
  xmin = -1.2,
  xmax =  1.2,
  ymin = -1.2,
  ymax =  1.2,
  zmin = -1.2,
  zmax =  1.2,
  nx   = 40,
  ny   = 40,
  nz   = 40,
  iso  = 0,
] ;
\stoptyping

Here the surface is where \type {code} evaluates to \type {iso}; in this
example, that value is \type {0}. The same sphere could be written as
\type {x*x + y*y + z*z} with \type {iso = 1}; subtracting \type {1} in the
expression and keeping \type {iso = 0} is just a convenient alternative. The
bounds define the box to search, and \type {nx}, \type {ny}, and \type {nz}
control its sampling. All three counts are required and must be larger than one.
The other settings handle naming, lighting normals, and material selection.

You can also supply the implicit function as a \LUA\ function through
\type {pfunction}. It receives \type {x}, \type {y}, and \type {z}, and returns
the scalar value. Return \type {false} when a sampled point lies outside the
part of the domain to draw. This lets you describe surfaces on specific,
nonrectangular domains.

The function is defined in the extended \type {xmath} table; inside a scene
expression the same table is reached as \type {math}, so the reference below
uses \type {pfunction = "math.paraboloid_domain"}.

\starttyping[option=LUA]
function xmath.paraboloid_domain(x,y,z)
    local r2 = x^2 + y^2
    if r2 < 0.25 or r2 > 1 then
        return false
    end
    return z - r2
end
\stoptyping

\startluacode
function xmath.paraboloid_domain(x,y,z)
    local r2 = x^2 + y^2
    if r2 < 0.25 or r2 > 1 then
        return false
    end
    return z - r2
end
\stopluacode

This uses that function to draw a paraboloid with a circular hole:

\startbuffer[implicit-lua-domain]
lmt_scene_start [
  width       = 6cm,
  crop        = true,
  resolution  = 100,
  supersample = 2,
  projection  = "perspective",
  light       = (0,-1,-2),
  eye         = (3,-2,3),
  intensity   = 1.2,
] ;

lmt_scene_implicit [
  pfunction = "math.paraboloid_domain",
  xmin      = -1.25,
  xmax      =  1.25,
  ymin      = -1.25,
  ymax      =  1.25,
  zmin      = -1.25,
  zmax      =  1.25,
  nx        = 120,
  ny        = 120,
  nz        = 120,
] ;

lmt_scene_stop ;
\stopbuffer

\typebuffer[implicit-lua-domain][option=MP]

\startplacefloat
  [figure]
  [title={An implicit paraboloid restricted by a \LUA-defined circular domain.}]
  \processMPbuffer[implicit-lua-domain]
\stopplacefloat

For an annular domain, add a second test such as
\type {or r2 > 1} to the condition. Because the predicate is sampled on the
implicit grid, increase \type {nx} and \type {ny} when a curved domain boundary
needs more detail.

The example uses approximately \type {120^3} samples. This is a practical
balance for a manual example: \type {200^3} gives a finer implicit grid but
requires substantially more time and memory, especially when the document is
compiled repeatedly (which is why we can cache).

Higher values of \type {nx}, \type {ny}, and \type {nz} add detail but can make
rendering much slower. A grid cell larger than a thin feature can miss it,
especially when the feature is small or disconnected. The finite bounds are a
hard clipping box, so geometry outside them is not extracted. Start with a box
that contains the part of the surface you need, then increase the counts until
the visible structure stops changing. Lower counts are useful for draft renders.
Unlike the graph and parametric commands, these counts describe sampling points,
not mesh intervals.

When analytic derivatives are available, give all three of \type {dfdx},
\type {dfdy}, and \type {dfdz}. They are used to make smoother and more accurate
lighting normals. They improve the normals on the sampled mesh, but do not
replace the mesh extraction or recover a component that the grid has missed.

The comparison below keeps the equation, bounds, and derivatives fixed and
changes only the grid resolution:

\startbuffer[implicit-resolution-coarse]
lmt_scene_start [
  width = 5.5cm,
  crop  = true,
] ;

lmt_scene_implicit [
  code = "x*x + y*y + z*z - 1",
  dfdx = "2*x",
  dfdy = "2*y",
  dfdz = "2*z",
  xmin = -1.2,
  xmax =  1.2,
  ymin = -1.2,
  ymax =  1.2,
  zmin = -1.2,
  zmax =  1.2,
  nx   = 10,
  ny   = 10,
  nz   = 10,
  iso  = 0,
] ;

lmt_scene_stop ;
\stopbuffer

\startbuffer[implicit-resolution-fine]
lmt_scene_start [
  width = 5.5cm,
  crop  = true,
] ;

lmt_scene_implicit [
  code = "x*x + y*y + z*z - 1",
  dfdx = "2*x",
  dfdy = "2*y",
  dfdz = "2*z",
  xmin = -1.2,
  xmax =  1.2,
  ymin = -1.2,
  ymax =  1.2,
  zmin = -1.2,
  zmax =  1.2,
  nx   = 32,
  ny   = 32,
  nz   = 32,
  iso  = 0,
] ;

lmt_scene_stop ;
\stopbuffer

\typebuffer[implicit-resolution-coarse][option=MP]

\startplacefloat
  [figure]
  [title={The same implicit sphere sampled on a coarse and a fine grid.}]
  \startcombination[nx=2,ny=1]
    {\processMPbuffer[implicit-resolution-coarse]}{coarse \type {nx = ny = nz = 10}}
    {\processMPbuffer[implicit-resolution-fine]}  {fine \type {nx = ny = nz = 32}}
  \stopcombination
\stopplacefloat

The coarse mesh shows the faceting caused by its larger cells. The fine mesh
follows the curved surface more closely, but costs more to sample.

The implicit-surface keywords are:

\startmixedcolumns

\keywordindex {id}
\keyword {id} gives the surface a name so that it can be referred to elsewhere,
for example by an intersection.

\keywordindex {code}
\keyword {code} gives \m {F(x,y,z)} as an expression in \type{x}, \type{y},
and \type{z}. The extracted surface is where this function reaches the chosen
\type {iso} value.

\keywordindex {pfunction}
\keyword {pfunction} gives a \LUA\ expression that evaluates to a function of
\type {x}, \type {y}, and \type {z}. It takes precedence over \type {code}; the
function can return \type {false} to omit a sampled point and restrict the
surface to a nonrectangular domain.

\keywordindex {splitdepth}
\keyword {splitdepth} sets the maximum number of times a cube is subdivided when
its corners contain a mixture of numeric and nonnumeric results from the point
function. This provides local refinement at a domain boundary made with
\type {pfunction}; an integer value of \type {1} subdivides a mixed cube once,
while larger values allow deeper refinement. The default is \type {2}, and
values outside the range \type {1}--\type {5} use that default. Cubes whose
corners are all numeric are polygonized immediately, while cubes with no
numeric corners are discarded, so this setting does not refine an ordinary
implicit function that returns a number everywhere.

\keywordindex {nfunction}
\keyword {nfunction} gives a \LUA\ expression that evaluates to a function returning
the normal components \type {(nx,ny,nz)} for \type {x}, \type {y}, and \type {z}.
It takes precedence over the three derivative expressions when normals are
calculated, while \type {code} or \type {pfunction} still defines the surface.

\keywordindex {dfdx}
\keyword {dfdx} gives the partial derivative with respect to \type{x}. When all
three derivatives are given, they are used for smoother and more accurate
normals.

\keywordindex {dfdy}
\keyword {dfdy} gives the partial derivative with respect to \type{y}. When all
three derivatives are given, they are used for smoother and more accurate
normals.

\keywordindex {dfdz}
\keyword {dfdz} gives the partial derivative with respect to \type{z}. When all
three derivatives are given, they are used for smoother and more accurate
normals.

\keywordindex {normal}
\keyword {normal} selects how normals are made for lighting. \type{smooth}
interpolates normals across the mesh, \type{flat} exposes the triangle faces,
\type{up} uses \type{(0,0,1)} everywhere, and \type{auto} estimates the
gradient of the implicit function from nearby samples.

\keywordindex {normalstep}
\keyword {normalstep} sets the finite-difference step used for numerical gradient
normals and for normal queries. A positive value overrides the automatic step;
otherwise the step is derived from the implicit bounding box.

\keywordindex {xmin}
\keyword {xmin} sets the lower bound of the sampled box in the \type{x}
direction.

\keywordindex {xmax}
\keyword {xmax} sets the upper bound of the sampled box in the \type{x}
direction.

\keywordindex {ymin}
\keyword {ymin} sets the lower bound of the sampled box in the \type{y}
direction.

\keywordindex {ymax}
\keyword {ymax} sets the upper bound of the sampled box in the \type{y}
direction.

\keywordindex {zmin}
\keyword {zmin} sets the lower bound of the sampled box in the \type{z}
direction.

\keywordindex {zmax}
\keyword {zmax} sets the upper bound of the sampled box in the \type{z}
direction.

\keywordindex {nx}
\keyword {nx} sets the number of sampling points in the \type{x} direction. It
is required and has to be larger than one. Higher values can capture more detail
but take more work.

\keywordindex {ny}
\keyword {ny} sets the number of sampling points in the \type{y} direction. It
is required and has to be larger than one. Higher values can capture more detail
but take more work.

\keywordindex {nz}
\keyword {nz} sets the number of sampling points in the \type{z} direction. It
is required and has to be larger than one. Higher values can capture more detail
but take more work.

\keywordindex {iso}
\keyword {iso} sets the function value that will be extracted as the surface.
The default sphere example uses \type {iso = 0}.

\keywordindex {material}
\keyword {material} selects the named material used when rendering the surface.

\stopmixedcolumns

\stopsection

\startsection[reference=parametric-curves,title=Parametric curves]

A parametric curve is a one-parameter version of a parametric surface. For each
sampled value of \m {t}, three expressions give a point \m {(x,y,z)} in space.
Use curves for paths, space curves, helper lines, and any curve that should use
the same depth handling as the surrounding surfaces.

The curve is sampled into points and rendered as a connected line in the scene
raster, so surfaces can hide parts behind them. Curves obtain their color from
their selected material, not from a separate curve-color setting; the material's
\type {dx} and \type {dy} values control the line-like thickness in raster pixels.

Add a parametric curve with

\starttyping[option=MP]
lmt_scene_curve [ ... ]
\stoptyping

A curve might look like this:

\starttyping[option=MP]
lmt_scene_curve [
  id   = "helix",
  x    = "cos(t)",
  y    = "sin(t)",
  z    = ".15*t",
  tmin = 0,
  tmax = 6*pi,
  nt   = 160,
] ;
\stoptyping

The coordinate expressions describe the path. The \type{t} bounds choose the
part of the parameter line to sample, and \type{nt} or \type{tstep} controls the
number of samples. More samples make a curved path smoother, but also add more
segments to the rendered scene.

The curve keywords are:

\startmixedcolumns

\keywordindex {id}
\keyword {id} gives the curve a name in the scene.

\keywordindex {x}
\keyword {x} gives the \type{x(t)} coordinate as an expression in \type{t}. It is
evaluated at each sampled parameter value.

\keywordindex {y}
\keyword {y} gives the \type{y(t)} coordinate as an expression in \type{t}. It is
evaluated at each sampled parameter value.

\keywordindex {z}
\keyword {z} gives the \type{z(t)} coordinate as an expression in \type{t}. It is
evaluated at each sampled parameter value.

\keywordindex {tmin}
\keyword {tmin} sets the lower bound of the sampled parameter interval.

\keywordindex {tmax}
\keyword {tmax} sets the upper bound of the sampled parameter interval.

\keywordindex {nt}
\keyword {nt} sets the number of sampling points along the curve. Both endpoints
are included, so the visible step size is based on \type{nt-1} intervals.

\keywordindex {tstep}
\keyword {tstep} requests an approximate sampling step in the parameter direction.
When it is positive, the point count is derived from the interval length and the
actual step is adjusted so that the curve still ends at \type{tmax}.

\keywordindex {material}
\keyword {material} selects the material used when rendering the curve. For curves,
the material controls the visible color, opacity, on-top behavior, and the raster
size settings used for point-like or line-like meshes.

\stopmixedcolumns

\startbuffer[curve-helix]
lmt_scene_start [
  width = 6cm,
  fov   = 90,
] ;

lmt_scene_curve [
  x    = "cos(t)",
  y    = "sin(t)",
  z    = ".15*t",
  tmin = 0,
  tmax = 6*pi,
  nt   = 300,
] ;

lmt_scene_stop ;
\stopbuffer

\startbuffer[curve-helix2]
lmt_scene_start [
  width = 6cm,
  fov   = 90,
] ;

lmt_scene_material [
    name      = "helix",
    opacity   = 1,
    dx        = 6,
    dy        = 6,
] ;

lmt_scene_curve [
  id   = "helix",
  x    = "cos(t)",
  y    = "sin(t)",
  z    = ".15*t",
  tmin = 0,
  tmax = 6*pi,
  nt   = 300,
  material = "helix",
] ;

lmt_scene_stop ;
\stopbuffer

\startplacefloat
  [figure]
  [title={A space curve with the default and a custom line material.}]
  \startcombination[nx=2,ny=1]
    {\startframed[TightBox]\processMPbuffer[curve-helix]\stopframed}{default}
    {\startframed[TightBox]\processMPbuffer[curve-helix2]\stopframed}{custom material}
  \stopcombination
\stopplacefloat

\startbuffer[curve-knot]
lmt_scene_start [
  width = 6cm,
  fov   = 90,
] ;

lmt_scene_material [
    name      = "knot",
    opacity   = 1,
    dx        = 12,
    dy        = 12,
] ;

lmt_scene_curve [
  x = "(2 + cos(3*t)) * cos(2*t)",
  y = "(2 + cos(3*t)) * sin(2*t)",
  z = "sin(3*t)",
  tmin = 0,
  tmax = 2*pi,
  nt   = 300,
  material = "knot",
] ;

lmt_scene_stop ;
\stopbuffer

\startbuffer[curve-sphericalwoven]
lmt_scene_start [
  width = 6cm,
  fov   = 90,
] ;

lmt_scene_material [
    name      = "knot",
    opacity   = 0.5,
    dx        = 4,
    dy        = 4,
] ;

lmt_scene_curve [
  x = "(1 + 0.22*cos(7*t))*sin(2*t)*cos(3*t)",
  y = "(1 + 0.22*cos(7*t))*sin(2*t)*sin(3*t)",
  z = "(1 + 0.22*cos(7*t))*cos(2*t)",
  tmin = 0,
  tmax = 2*pi,
  nt   = 300,
  material = "knot",
] ;

lmt_scene_stop ;
\stopbuffer

\startplacefloat
  [figure]
  [title={Two sampled space curves.}]
  \startcombination[nx=2,ny=1]
    {\startframed[TightBox]\processMPbuffer[curve-knot]\stopframed}{knot}
    {\startframed[TightBox]\processMPbuffer[curve-sphericalwoven]\stopframed}{woven curve}
  \stopcombination
\stopplacefloat

\stopsection

\startsection[reference=intersections,title=Intersections]

An intersection finds where two objects already in the scene meet. The result is
added as point-like or segment-like geometry and can be rendered with the source
objects. Use intersections for level curves, cutting planes, and the places where
two sampled surfaces pass through each other.

Intersections are computed from sampled geometry, so the result is a sampled
visualization rather than an exact geometric operation. Its appearance depends
on the source resolution and the selected method; finer source meshes usually
give a cleaner result. The \type {tolerance} is a small numerical threshold used
when crossings are classified and nearby results are joined; the useful value
depends on the scale of the objects. The most dependable case is a generated
surface or mesh together with a finite plane.
Intersections between two generated surfaces are still experimental and may
produce no result for some object types.

Request an intersection with

\starttyping[option=MP]
lmt_scene_intersection [ ... ]
\stoptyping

An intersection might look like this:

\starttyping[option=MP]
lmt_scene_intersection [
  id        = "slice",
  first     = "surface",
  second    = "cutplane",
  tolerance = .001,
  material  = "default",
] ;
\stoptyping

Here \type{surface} and \type{cutplane} are the \type{id} values of objects that
have already been added to the same scene. A graph or mesh can be intersected
with a finite plane patch. Mesh--mesh requests are experimental; for a cutting
curve, start with the plane form shown here.

The intersection keywords are:

\startmixedcolumns

\keywordindex {id}
\keyword {id} gives the intersection request a name. The important names for
this lookup are the source object names given by \type{first} and
\type{second}.

\keywordindex {first}
\keyword {first} selects the first object in the intersection. It should match
the \type{id} of an earlier object in the same scene.

\keywordindex {second}
\keyword {second} selects the second object in the intersection. It should match
the \type{id} of another earlier object in the same scene.

\keywordindex {tolerance}
\keyword {tolerance} sets the numerical threshold used by the selected
intersection method. It should normally remain small relative to the coordinate
scale of the two objects.

\keywordindex {material}
\keyword {material} selects the material used when rendering the generated
intersection result. For these line-like results, the material controls the
visible color, opacity, on-top behavior, and raster size settings.

\stopmixedcolumns

\index {opacity}%
The complete example below uses a parametric torus and a finite horizontal
plane. By the time the intersection is calculated, the torus is a triangle mesh,
so this exercises the mesh--plane path. The red loops are the generated
line-like intersection. The plane is translucent so that both the mesh and the
result remain visible.

\startbuffer[intersection-mesh-plane]
lmt_scene_start [
  width      = 6cm,
  crop       = true,
  eye        = (3,-4,2.6),
  target     = (0,0,0),
  light      = (-1,-1,-2),
] ;

lmt_scene_material [
  name      = "torus",
  diffuse   = (.20,.48,.82),
  ambient   = .25,
  opacity   = 1,
] ;

lmt_scene_material [
  name      = "slice",
  diffuse   = (.72,.78,.68),
  opacity   = .20,
  twosided  = true,
] ;

lmt_scene_material [
  name      = "section",
  diffuse   = (.95,.05,.03),
  dx        = 6,
  dy        = 6,
  ontop     = true,
] ;

lmt_scene_parametric [
  id       = "torus",
  x        = "(1+.35*cos(v))*cos(u)",
  y        = "(1+.35*cos(v))*sin(u)",
  z        = ".35*sin(v)",
  umin     = 0,
  umax     = 2*pi,
  vmin     = 0,
  vmax     = 2*pi,
  nu       = 48,
  nv       = 24,
  material = "torus",
] ;

lmt_scene_plane [
  id       = "slice",
  normal   = (0,0,1),
  center   = (0,0,0),
  usize    = 3,
  vsize    = 3,
  material = "slice",
] ;

lmt_scene_intersection [
  id        = "section",
  first     = "torus",
  second    = "slice",
  tolerance = .001,
  material  = "section",
] ;

lmt_scene_stop ;
\stopbuffer

\typebuffer[intersection-mesh-plane][option=MP]

\startplacefloat
  [figure]
  [title={A tested mesh-plane intersection of a torus and a finite plane.}]
  \processMPbuffer[intersection-mesh-plane]
\stopplacefloat

\stopsection

\startsection[reference=reference-geometry,title=Reference geometry]

Reference geometry provides finite planes, segments, grids, axes, and ticks that
can be mixed with sampled surfaces and meshes. These helpers use raster-pixel
widths where they draw lines, so their visual weight is independent of the
document's physical dimensions.
\index {depth testing}%
\index {line widths}%
The line-like helpers are depth-tested with the
other scene geometry; grids, axes, and ticks additionally describe an origin,
direction vectors, bounds or positions, color, and width.

\startsubsubject[reference=planes,title=Planes]

A plane object is a flat rectangular patch in 3D space. Use it as a visible
reference sheet, a transparent cutting plane, or a simple surface for comparison.
The plane is positioned by a point and oriented by a normal direction; its two
sizes set the finite visible patch.

The patch uses the same camera, light, material, and depth handling as the other
scene objects. Its local \type{u} and \type{v} directions are chosen
automatically from the normal.

Add a plane with

\starttyping[option=MP]
lmt_scene_plane [ ... ]
\stoptyping

A plane might look like this:

\starttyping[option=MP]
lmt_scene_plane [
  id       = "slice",
  normal   = (0,0,1),
  center   = (0,0,.35),
  usize    = 3,
  vsize    = 3,
  material = "default",
] ;
\stoptyping

The patch is centered at \type{(0,0,.35)} and has the normal \type{(0,0,1)}.
The two sizes set the visible rectangle in directions perpendicular to that
normal. For a helper or cutting surface, a transparent or on-top material often
makes the plane easier to read with the rest of the scene.

The plane keywords are:

\startmixedcolumns

\keywordindex {id}
\keyword {id} gives the plane a name so that it can be referred to elsewhere,
for example by an intersection.

\keywordindex {normal}
\keyword {normal} sets the plane normal. It gives the orientation of the plane
and is normalized internally, so it does not have to be a unit vector.

\keywordindex {center}
\keyword {center} sets the center point of the finite plane patch.

\keywordindex {usize}
\keyword {usize} sets the size of the plane patch in the local \type{u}
direction. This direction lies in the plane and is chosen automatically from the
normal.

\keywordindex {vsize}
\keyword {vsize} sets the size of the plane patch in the local \type{v}
direction. Together with \type{usize} it determines the visible rectangle.

\keywordindex {material}
\keyword {material} selects the material used when rendering the plane. This is
where color, transparency, and on-top behavior are normally controlled.

\stopmixedcolumns

Here is a standalone plane scene:

\startbuffer[plane-standalone]
lmt_scene_start [
  width      = 6cm,
  resolution = 120,
  crop       = true,
  eye        = (3,-4,2.6),
  target     = (0,0,0),
] ;

lmt_scene_material [
  name      = "plane",
  diffuse   = (.25,.55,.85),
  twosided  = true,
  opacity   = .75,
] ;

lmt_scene_plane [
  id       = "plane",
  normal   = (0,0,1),
  center   = (0,0,0),
  usize    = 3,
  vsize    = 2,
  material = "plane",
] ;

lmt_scene_stop ;
\stopbuffer

\typebuffer[plane-standalone][option=MP]

\startplacefloat
  [figure]
  [title={A standalone finite plane patch.}]
  \processMPbuffer[plane-standalone]
\stopplacefloat

\stopsubsubject

\startsubsubject[reference=segments,title=Segments]

A segment is a straight line between two 3D points. Use it for explicit helper
lines: marking a distance, showing a direction, or highlighting an edge that is
not part of a surface.

Segments are rendered as line-like scene geometry, so they take part in the same
depth handling as curves and intersections. A segment can therefore disappear
behind an opaque surface unless it is styled through the usual on-top mechanisms.

Add a visible 3D segment with

\starttyping[option=MP]
lmt_scene_segment [ ... ]
\stoptyping

A segment might look like this:

\starttyping[option=MP]
lmt_scene_segment [
  id     = "baseline",
  first  = (-1,0,0),
  second = ( 1,0,0),
  color  = red,
  width  = 1,
] ;
\stoptyping

The segment keywords are:

\startmixedcolumns

\keywordindex {id}
\keyword {id} gives the segment helper a scene name.

\keywordindex {first}
\keyword {first} sets the first endpoint as a 3D point.

\keywordindex {second}
\keyword {second} sets the second endpoint as a 3D point.

\keywordindex {color}
\keyword {color} sets the segment color.

\keywordindex {width}
\keyword {width} sets the rendered segment width in raster pixels. This helper
width controls the raster thickness directly, rather than the document width.

\stopmixedcolumns

Here is a standalone segment scene:

\startbuffer[segment-standalone]
lmt_scene_start [
  width      = 6cm,
  resolution = 120,
  crop       = true,
  eye        = (3,-4,2.6),
  target     = (0,0,.3),
] ;

lmt_scene_segment [
  id     = "segment",
  first  = (-1.4,-.5,0),
  second = ( 1.4,.8,1),
  color  = red,
  width  = 4,
] ;

lmt_scene_stop ;
\stopbuffer

\typebuffer[segment-standalone][option=MP]

\startplacefloat
  [figure]
  [title={A standalone 3D segment with a four-pixel raster width.}]
  \processMPbuffer[segment-standalone]
\stopplacefloat

\stopsubsubject

\startsubsubject[reference=grids,title=Grids]

A grid is a set of straight helper lines in a plane. It is defined by an origin
and two direction vectors, \type{uaxis} and \type{vaxis}. The bounds and step
sizes then say where the grid lines are placed in those two local directions.

They work well as reference planes under or through a surface, or as a lightweight
coordinate frame when a full axis setup would be too much.

Add a grid with

\starttyping[option=MP]
lmt_scene_grid [ ... ]
\stoptyping

A grid might look like this:

\starttyping[option=MP]
lmt_scene_grid [
  origin = (0,0,0),
  uaxis  = (1,0,0),
  vaxis  = (0,1,0),
  umin   = -2,
  umax   =  2,
  ustep  = .5,
  vmin   = -2,
  vmax   =  2,
  vstep  = .5,
  color  = (.62,.62,.62),
  width  = 1,
] ;
\stoptyping

The grid keywords are:

\startmixedcolumns

\keywordindex {id}
\keyword {id} gives the grid helper a scene name.

\keywordindex {origin}
\keyword {origin} sets the point where the local \type{u} and \type{v}
coordinates start from.

\keywordindex {uaxis}
\keyword {uaxis} sets the first direction vector of the grid. Its length also
sets the scale of one unit in the local \type{u} direction.

\keywordindex {vaxis}
\keyword {vaxis} sets the second direction vector of the grid. Its length also
sets the scale of one unit in the local \type{v} direction.

\keywordindex {umin}
\keyword {umin} sets the lower bound in the local \type{u} direction.

\keywordindex {umax}
\keyword {umax} sets the upper bound in the local \type{u} direction.

\keywordindex {ustep}
\keyword {ustep} sets the spacing between grid lines in the local \type{u}
direction.

\keywordindex {vmin}
\keyword {vmin} sets the lower bound in the local \type{v} direction.

\keywordindex {vmax}
\keyword {vmax} sets the upper bound in the local \type{v} direction.

\keywordindex {vstep}
\keyword {vstep} sets the spacing between grid lines in the local \type{v}
direction.

\keywordindex {color}
\keyword {color} sets the color of the visible grid.

\keywordindex {width}
\keyword {width} sets the line width of the visible grid in raster pixels.

\stopmixedcolumns

The examples show the grid by itself and with an opaque surface. In the second
case it is ordinary scene geometry, so the surface hides the parts behind it.
Both examples use the grid from the code above.

\startbuffer[grid-only]
lmt_scene_start [
  width      = 5cm,
  crop       = true,
  eye        = (3,-4,2.6),
] ;

lmt_scene_grid [
  origin = (0,0,0),
  uaxis  = (1,0,0),
  vaxis  = (0,1,0),
  umin   = -2,
  umax   =  2,
  ustep  = .5,
  vmin   = -2,
  vmax   =  2,
  vstep  = .5,
  color  = (.62,.62,.62),
  width  = 1,
] ;

lmt_scene_stop ;
\stopbuffer

\startbuffer[grid-surface]
lmt_scene_start [
  width      = 5cm,
  crop       = true,
  eye        = (3,-4,2.6),
] ;

lmt_scene_material [
  name      = "gridhill",
  diffuse   = (.62,.72,.86),
  specular  = (.1,.1,.1),
  shininess = 12,
  ambient   = .25,
] ;

lmt_scene_surface [
  code     = ".7*exp(-.6*(x*x+y*y)) - 0.2",
  dfdx     = "-.84*x*exp(-.6*(x*x+y*y))",
  dfdy     = "-.84*y*exp(-.6*(x*x+y*y))",
  xmin     = -1.8,
  xmax     =  1.8,
  ymin     = -1.8,
  ymax     =  1.8,
  nx       = 48,
  ny       = 48,
  material = "gridhill",
] ;

lmt_scene_grid [
  origin = (0,0,0),
  uaxis  = (1,0,0),
  vaxis  = (0,1,0),
  umin   = -2,
  umax   =  2,
  ustep  = .5,
  vmin   = -2,
  vmax   =  2,
  vstep  = .5,
  color  = (.62,.62,.62),
  width  = 1,
] ;

lmt_scene_stop ;
\stopbuffer

\startplacefloat
  [figure]
  [title={A grid on its own and depth-tested with a surface.}]
  \startcombination[nx=2,ny=1]
    {\processMPbuffer[grid-only]}   {grid only}
    {\processMPbuffer[grid-surface]}{with surface}
  \stopcombination
\stopplacefloat

\stopsubsubject

\startsubsubject[reference=axes,title=Axes]

An axis overlay draws the three coordinate axes as straight 3D segments. It is a
compact way to show orientation and scale without drawing a full grid. The axes
share one origin, but each direction vector can be changed, so the overlay can
also be used for a local coordinate frame.

Add an axis with

\starttyping[option=MP]
lmt_scene_axis [ ... ]
\stoptyping

An axis overlay might look like this:

\starttyping[option=MP]
lmt_scene_axis [
  origin = (0,0,0),
  xaxis  = (1,0,0),
  yaxis  = (0,1,0),
  zaxis  = (0,0,1),
  xmin   = -2,
  xmax   =  2,
  ymin   = -2,
  ymax   =  2,
  zmin   =  0,
  zmax   =  1,
  color  = black,
  width  = 1,
] ;
\stoptyping

The axis keywords are:

\startmixedcolumns

\keywordindex {id}
\keyword {id} gives the axis helper a scene name.

\keywordindex {origin}
\keyword {origin} sets the common origin point of the three axes.

\keywordindex {xaxis}
\keyword {xaxis} sets the direction vector of the local \type{x} axis. Its length
sets the scale of one local \type{x} unit.

\keywordindex {yaxis}
\keyword {yaxis} sets the direction vector of the local \type{y} axis. Its length
sets the scale of one local \type{y} unit.

\keywordindex {zaxis}
\keyword {zaxis} sets the direction vector of the local \type{z} axis. Its length
sets the scale of one local \type{z} unit.

\keywordindex {xmin}
\keyword {xmin} sets the lower bound of the visible local \type{x} axis.

\keywordindex {xmax}
\keyword {xmax} sets the upper bound of the visible local \type{x} axis.

\keywordindex {ymin}
\keyword {ymin} sets the lower bound of the visible local \type{y} axis.

\keywordindex {ymax}
\keyword {ymax} sets the upper bound of the visible local \type{y} axis.

\keywordindex {zmin}
\keyword {zmin} sets the lower bound of the visible local \type{z} axis.

\keywordindex {zmax}
\keyword {zmax} sets the upper bound of the visible local \type{z} axis.

\keywordindex {color}
\keyword {color} sets the color of the visible axes.

\keywordindex {width}
\keyword {width} sets the line width of the visible axes in raster pixels.

\stopmixedcolumns


\startbuffer[axis-surface]
lmt_scene_start [
  width      = 5cm,
  crop       = true,
  eye        = (3,-4,2.6),
] ;

lmt_scene_material [
  name      = "axishill",
  diffuse   = (.62,.72,.86),
  specular  = (.1,.1,.1),
  shininess = 12,
  ambient   = .25,
] ;

lmt_scene_surface [
  code     = ".7*exp(-.6*(x*x+y*y)) - 0.2",
  dfdx     = "-.84*x*exp(-.6*(x*x+y*y))",
  dfdy     = "-.84*y*exp(-.6*(x*x+y*y))",
  xmin     = -1.8,
  xmax     =  1.8,
  ymin     = -1.8,
  ymax     =  1.8,
  nx       = 48,
  ny       = 48,
  material = "axishill",
] ;

lmt_scene_axis [
  origin = (0,0,0),
  xaxis  = (1,0,0),
  yaxis  = (0,1,0),
  zaxis  = (0,0,1),
  xmin   = -2,
  xmax   =  2,
  ymin   = -2,
  ymax   =  2,
  zmin   =  0,
  zmax   =  1,
  color  = black,
  width  = 1,
] ;

lmt_scene_stop ;
\stopbuffer

\startplacefloat
  [figure]
  [title={Axes depth-tested together with a surface.}]
  \processMPbuffer[axis-surface]
\stopplacefloat

\stopsubsubject

\startsubsubject[reference=ticks,title=Ticks]

Ticks are short helper marks placed along one direction. They are normally used
together with an axis: \type{axis} says where the tick positions lie, and
\type{tickaxis} says which way the small marks extend.

The tick command only draws the marks themselves; labels, when needed, are added
separately.

Add ticks with

\starttyping[option=MP]
lmt_scene_ticks [ ... ]
\stoptyping

A tick setup might look like this:

\starttyping[option=MP]
lmt_scene_ticks [
  origin   = (0,0,0),
  axis     = (1,0,0),
  tickaxis = (0,1,0),
  first    = -2,
  last     =  2,
  step     =  1,
  size     = .15,
  color    = red,
  width    = 1,
] ;
\stoptyping

The tick keywords are:

\startmixedcolumns

\keywordindex {id}
\keyword {id} gives the tick helper a scene name.

\keywordindex {origin}
\keyword {origin} sets the point from which tick positions are measured.

\keywordindex {axis}
\keyword {axis} sets the direction vector along which the ticks are placed. Its
length sets the scale of one tick-position unit.

\keywordindex {tickaxis}
\keyword {tickaxis} sets the direction vector along which each tick mark extends.

\keywordindex {first}
\keyword {first} sets the first tick position on the axis.

\keywordindex {last}
\keyword {last} sets the last tick position on the axis.

\keywordindex {step}
\keyword {step} sets the spacing between consecutive tick positions.

\keywordindex {size}
\keyword {size} sets the length of each tick mark, measured along
\type{tickaxis}.

\keywordindex {color}
\keyword {color} sets the color of the visible ticks.

\keywordindex {width}
\keyword {width} sets the line width of the visible ticks in raster pixels.

\stopmixedcolumns


\startbuffer[ticks-surface]
lmt_scene_start [
  width       = 5cm,
%   supersample = 2,
  crop        = true,
  eye         = (3,-4,2.6),
] ;

lmt_scene_material [
  name      = "axishill",
  diffuse   = (.62,.72,.86),
  specular  = (.1,.1,.1),
  shininess = 12,
  ambient   = .25,
] ;

lmt_scene_surface [
  code     = ".7*exp(-.6*(x*x+y*y)) - 0.2",
  dfdx     = "-.84*x*exp(-.6*(x*x+y*y))",
  dfdy     = "-.84*y*exp(-.6*(x*x+y*y))",
  xmin     = -1.8,
  xmax     =  1.8,
  ymin     = -1.8,
  ymax     =  1.8,
  nx       = 48,
  ny       = 48,
  material = "axishill",
] ;

lmt_scene_axis [
  origin = (0,0,0),
  xaxis  = (1,0,0),
  yaxis  = (0,1,0),
  zaxis  = (0,0,1),
  xmin   = -2,
  xmax   =  2,
  ymin   = -2,
  ymax   =  2,
  zmin   =  0,
  zmax   =  1,
  color  = black,
  width  = 1,
] ;

lmt_scene_ticks [
  origin   = (0,0,0),
  axis     = (1,0,0),
  tickaxis = (0,1,0),
  first    = -2,
  last     =  2,
  step     =  1,
  size     = .15,
  color    = red,
  width    = 1,
] ;

lmt_scene_stop ;
\stopbuffer

\startplacefloat
  [figure]
  [title={Axes and ticks depth-tested together with a surface.}]
  \processMPbuffer[ticks-surface]
\stopplacefloat

\stopsubsubject

\stopsection

\startsection[reference=complete-scene,title=Putting it all together]

The complete example below brings the basic workflow together. It combines a
named surface, three named materials, a finite plane, their sampled
intersection, and coordinate axes. The intersection refers to the two objects
by their \type {id} values.

\startbuffer[scene-first]
lmt_scene_start [
  width      = 6cm,
  resolution = 120,
  crop       = true,
  eye        = (3.2,-4.4,3.0),
  target     = (0,0,.1),
] ;

lmt_scene_material [
  name      = "hill",
  diffuse   = (.25,.55,.85),
  specular  = (.12,.12,.12),
  shininess = 12,
  ambient   = .22,
] ;

lmt_scene_material [
  name      = "cut",
  diffuse   = (1,.65,.08),
  opacity   = .28,
  twosided  = true,
] ;

lmt_scene_material [
  name      = "section",
  diffuse   = (.9,.08,.06),
  dx        = 5,
  dy        = 5,
  ontop     = true,
] ;

lmt_scene_surface [
  id       = "hill",
  code     = ".8*exp(-.35*(x*x+y*y)) - .25",
  dfdx     = "-.56*x*exp(-.35*(x*x+y*y))",
  dfdy     = "-.56*y*exp(-.35*(x*x+y*y))",
  xmin     = -2,
  xmax     =  2,
  ymin     = -2,
  ymax     =  2,
  nx       = 56,
  ny       = 56,
  material = "hill",
] ;

lmt_scene_plane [
  id       = "cutplane",
  normal   = (0,0,1),
  center   = (0,0,.15),
  usize    = 3.6,
  vsize    = 3.6,
  material = "cut",
] ;

lmt_scene_intersection [
  id        = "section",
  first     = "hill",
  second    = "cutplane",
  tolerance = .01,
  material  = "section",
] ;

lmt_scene_axis [
  origin = (0,0,0),
  xaxis  = (1,0,0),
  yaxis  = (0,1,0),
  zaxis  = (0,0,1),
  xmin   = -2,
  xmax   =  2,
  ymin   = -2,
  ymax   =  2,
  zmin   = -.4,
  zmax   =  .8,
  color  = black,
  width  = 1,
] ;

lmt_scene_stop ;
\stopbuffer

\typebuffer[scene-first][option=MP]

\startplacefloat
  [figure]
  [title={A complete scene: a surface, a cutting plane, their intersection,
   and coordinate axes.}]
  \processMPbuffer[scene-first]
\stopplacefloat

Add the surface and plane before the intersection: \type {first} and
\type {second} refer to objects that already exist in the scene. The axis has no
such dependency, so it comes after the intersection to make the construction
order clear. All four objects use the same camera, lighting, materials, and
depth handling.

The intersection is a sampled graph--plane result for visual cutting curves.
Its numerical and tolerance-related limits are described in the
\goto {Intersections} [intersections] section. Use this scene as a workflow
example; the individual sections remain the command reference.

The \type {width} above is the displayed document size. Since no explicit
\type {bytewidth} or \type {byteheight} is given, \type {resolution} determines
the raster size used before the result is placed at that width. With
\type {crop = true}, the displayed bitmap is fitted to the projected bounds of
the complete scene, including the plane, intersection, and axes.

\type {lmt_scene_stop} renders the scene when no earlier render was requested. If
ordinary \METAPOST\ labels or markers need the rendered bounds or projected
positions, call \type {lmt_scene_render} after all scene objects and before
\type {lmt_scene_stop}; the query helpers are described in the
\goto {advanced annotation section} [annotation-helpers].

\stopsection

\startsection[reference=annotation-helpers,title=Rendering and annotation helpers (advanced)]

These helpers work on a rendered scene. Use them when ordinary \METAPOST\
material must be positioned relative to the result. Otherwise,
\type {lmt_scene_stop} renders and finishes the scene for you.

\type {lmt_scene_render} renders the scene once. After rendering,
\index {bytemap}%
\type {lmt_scene_width} and \type {lmt_scene_height} give the pixel dimensions of
the final raster. They need not equal the document dimensions requested with
\type {width} and \type {height}. The \type {lmt_scene_bytemap} and
\type {lmt_scene_bounds} helpers give access to the selected bytemap and its
displayed bounds.

The projection helper \type {lmt_scene_project_point} maps a 3D point to the
picture. The visibility helper \type {lmt_scene_project_visible} reports whether
that point is visible at its projected position. Their shorter aliases are
\type {lmt_scene_point} and \type {lmt_scene_visible}.

For an object whose \type {id} is known, \type {lmt_scene_calculate_normal}
returns a normal at a coordinate or parameter value; its shorter alias is
\type {lmt_scene_normal}. Make these queries after \type {lmt_scene_render} and
before \type {lmt_scene_stop}.

Here is a complete query-and-annotation example. The scene is rendered first;
the rest is ordinary \METAPOST\ code using the returned bounds, projected point,
visibility result, raster dimensions, and normal.

\startbuffer[scene-helpers]
lmt_scene_start [
  width      = 6cm,
  resolution = 120,
  crop       = true,
  eye        = (3,-4,2.6),
] ;

lmt_scene_parametric [
  id       = "surface",
  x        = "u",
  y        = "v",
  z        = ".7*exp(-.6*(u*u+v*v)) - .2",
  xu       = "1",
  yu       = "0",
  zu       = "-.84*u*exp(-.6*(u*u+v*v))",
  xv       = "0",
  yv       = "1",
  zv       = "-.84*v*exp(-.6*(u*u+v*v))",
  umin     = -1.8,
  umax     =  1.8,
  vmin     = -1.8,
  vmax     =  1.8,
  nu       = 48,
  nv       = 48,
] ;

lmt_scene_render ;

path     helper_bounds ;
pair     helper_probe ;
pair     helper_label ;
triplet  helper_normal ;

helper_bounds  := lmt_scene_bounds ;
helper_probe   := lmt_scene_point((0,0,-.5)) ;
helper_normal  := lmt_scene_normal("surface",(0,0)) ;
helper_label   := llcorner helper_bounds + (4,6) ;

draw helper_bounds withpen pencircle scaled .5 withcolor (.7,.7,.7) ;
draw fullcircle scaled 6 shifted helper_probe withcolor red ;
draw scenepoint "surface" (0,0,-.5)
  withpen pencircle scaled 2
  withcolor blue ;
draw scenenormal "surface" (0,0)
  withpen pencircle scaled .5
  withcolor green ;
draw textext("projected hidden point") shifted helper_probe shifted (4,16) ;

if not lmt_scene_visible((0,0,-.5)):
    draw textext("depth: hidden") shifted helper_probe shifted (4,-4) ;
fi ;

draw textext("raster " & decimal lmt_scene_width & " x " &
    decimal lmt_scene_height) shifted (helper_label shifted (0,14)) ;

draw textext("normal " & decimal redpart helper_normal & ", " &
    decimal greenpart helper_normal & ", " &
    decimal bluepart helper_normal) shifted (helper_label shifted (0,28)) ;

lmt_scene_stop ;
\stopbuffer

\typebuffer[scene-helpers][option=MP]

\startplacefloat
  [figure]
  [title={Using the rendered scene as a basis for ordinary \METAPOST\ annotations.}]
  \processMPbuffer[scene-helpers]
\stopplacefloat

The point returned by \type {lmt_scene_point} is already in the displayed
picture's coordinate system, so it can be used directly for a marker or label.
The example uses an analytic parametric surface, so the pair passed to
\type {lmt_scene_normal} is its \type {(u,v)} parameter pair. For an implicit
surface, pass a 3D coordinate instead. The normal itself is a \type {triplet};
the example formats its three components into a label. The projected hidden
point demonstrates that a point can still have a position even when
\type {lmt_scene_visible} reports that the depth test hides it.

\type {scenepoint} is a convenience form of \type {lmt_scene_point}: the object
tag comes first and the point follows it, and the result is an ordinary
projected \METAPOST\ pair. \type {scenenormal} uses the same tag-and-parameter
syntax but returns a small arrow picture. Both forms require
\type {lmt_scene_render} first; they are useful when the result is drawn
directly rather than stored for further calculations.

\stopsection

\startsection[reference=scene-cache,title=Scene cache (optional)]

You can cache a complete scene as a rendered bytemap. This helps when repeated
runs spend time generating the same geometry. Give the scene a stable
\type {name} and set \type {cache = true}:

\starttyping[option=MP]
lmt_scene_start [
  name       = "cached-plane",
  cache      = true,
  width      = 4cm,
  resolution = 120,
] ;

lmt_scene_plane [
  id     = "ground",
  normal = (0,0,1),
  center = (0,0,0),
  usize  = 3,
  vsize  = 3,
] ;

lmt_scene_stop ;
\stoptyping

The first run renders the scene and writes a cache file in the job directory.
A later run with the same job name, scene name, and scene description reuses
the cached bytemap instead of generating the 3D scene again. Changes to the
documented scene or object settings invalidate the entry and cause a fresh
render. Set \type {cache = false}, or omit \type {cache}, to bypass caching.

The cache stores the complete rendered scene, not individual meshes. Use a
stable, distinct \type {name} for each scene definition that should be cached.
Cache files are not removed automatically, so obsolete entries can be deleted
from the job directory. Caching requires an enabled \type {bytemap}, which is
the default; \type {bytemap = false} is not supported with \type {cache = true}.
You can use \typ {context --purge} to get rid of cache files.

Scenes with a \type {process} or \type {stipple} postprocessor are rendered
again when needed; those postprocessing passes are not reused from the scene
cache.

\stopsection

\startsection[reference=stippling,title=Stippling (advanced)]

Stippling converts the rendered scene into a field of colored dots. It is a
raster postprocessing step: geometry, camera, lighting, and materials are
resolved first, then the raster colors are converted to ink coverage. Depth
and surface normals help preserve silhouettes, creases, and separation between
visible surfaces.

Enable stippling with a table on \type {lmt_scene_start}. The example below
renders the same graph as a continuous raster and as a four-level stippled image:

\startbuffer[stipple-plain]
lmt_scene_start [
  width      = 5.5cm,
  resolution = 120,
  crop       = true,
  eye        = (3,-4,2.6),
  target     = (0,0,0),
  light      = (-1,-1,-2),
] ;

lmt_scene_surface [
  code = ".55*exp(-.07*(x*x+y*y))*(cos(2.2*x)+sin(2.0*y))",
  dfdx = ".55*exp(-.07*(x*x+y*y))*(-.14*x*(cos(2.2*x)+sin(2.0*y))-2.2*sin(2.2*x))",
  dfdy = ".55*exp(-.07*(x*x+y*y))*(-.14*y*(cos(2.2*x)+sin(2.0*y))+2.0*cos(2.0*y))",
  xmin = -3,
  xmax =  3,
  ymin = -3,
  ymax =  3,
  nx   = 64,
  ny   = 64,
] ;

lmt_scene_stop ;
\stopbuffer

\startbuffer[stipple-dotted]
lmt_scene_start [
  width      = 5.5cm,
  resolution = 120,
  crop       = true,
  eye        = (3,-4,2.6),
  target     = (0,0,0),
  light      = (-1,-1,-2),
  stipple    = [
    color           = (.15,.25,.55),
    backgroundcolor = (1,1,1),
    edgeboost       = .20,
    coveragegamma   = 1.25,
    levels          = 4,
    tonedots        = true,
    outline         = true,
  ],
] ;

lmt_scene_surface [
  code = ".55*exp(-.07*(x*x+y*y))*(cos(2.2*x)+sin(2.0*y))",
  dfdx = ".55*exp(-.07*(x*x+y*y))*(-.14*x*(cos(2.2*x)+sin(2.0*y))-2.2*sin(2.2*x))",
  dfdy = ".55*exp(-.07*(x*x+y*y))*(-.14*y*(cos(2.2*x)+sin(2.0*y))+2.0*cos(2.0*y))",
  xmin = -3,
  xmax =  3,
  ymin = -3,
  ymax =  3,
  nx   = 64,
  ny   = 64,
] ;

lmt_scene_stop ;
\stopbuffer

\typebuffer[stipple-dotted][option=MP]

\startplacefloat
  [figure]
  [title={The same graph rendered as a continuous raster and as a stippled image.}]
  \startcombination[nx=2,ny=1]
    {\processMPbuffer[stipple-plain]} {continuous raster}
    {\processMPbuffer[stipple-dotted]}{four-level stipple}
  \stopcombination
\stopplacefloat

The left image is the ordinary rendered raster. The right applies stippling
after lighting and visibility have been resolved. Stippling changes the raster,
not the scene geometry. Its appearance depends on raster resolution; use
\type {resolution} or explicit \type {bytewidth} and \type {byteheight} to
control the dot spacing. The default \type {dotradius = 0} uses one-pixel dots.
Larger radii can overlap and may need to be adjusted with the raster size.

The stipple settings are in the nested table passed to \type {stipple}:

\startmixedcolumns

\keywordindex {color}
\keyword {color} sets the ink color used for dots.

\keywordindex {backgroundcolor}
\keyword {backgroundcolor} sets the paper or background color. It is also used
for the intermediate colors produced by \type {tonedots}.

\keywordindex {nobackground}
\keyword {nobackground} keeps the existing raster background instead of
clearing it to the background color.

\keywordindex {levels}
\keyword {levels} selects evenly spaced tone levels. Give an integer from
\type {3} through \type {16}; omit it for binary stippling.

\keywordindex {tonedots}
\keyword {tonedots} uses colors between the background and ink colors for
intermediate tone levels. The default is true.

\keywordindex {errorthreshold}
\keyword {errorthreshold} sets the threshold used by binary stippling. The
default is \type {.5}; it has no effect when \type {levels} is active.

\keywordindex {coveragegamma}
\keyword {coveragegamma} reshapes the coverage curve before quantization. The
default is \type {1.25}; \type {1} leaves the curve unchanged.

\keywordindex {minimumcoverage}
\keywordindex {maximumcoverage}
\keyword {minimumcoverage/maximumcoverage} clamp coverage after edge boosting
and before the gamma adjustment. Their defaults are \type {0} and \type {1}.

\keywordindex {edgeboost}
\keyword {edgeboost} adds ink coverage near detected depth or normal edges. The
default is \type {.20}.

\keywordindex {usenormal}
\keyword {usenormal} includes surface-normal continuity in edge detection and
error diffusion. Set it to false to use depth information only.

\keywordindex {creaseangle}
\keyword {creaseangle} sets the normal-angle threshold for crease edges. The
default is \type {25} degrees.

\keywordindex {sameangle}
\keyword {sameangle} sets the normal-angle threshold used to decide whether
diffusion can cross neighboring pixels. The default is \type {15} degrees.

\keywordindex {dzsamemultiplier}
\keywordindex {dzedgemultiplier}
\keyword {dzsamemultiplier/dzedgemultiplier} scale the automatically estimated
depth differences for same-surface and edge tests. Their defaults are \type {2}
and \type {8}.

\keywordindex {mindepthstep}
\keyword {mindepthstep} supplies the minimum absolute depth step used by those
tests. The default is \type {1e-12}.

\keywordindex {serpentine}
\keyword {serpentine} alternates the scan direction on successive rows. The
default is true.

\keywordindex {dotradius}
\keyword {dotradius} sets the radius of filled dots in raster pixels. Zero
produces single-pixel dots.

\keywordindex {clipdots}
\keyword {clipdots} prevents larger dots from being painted onto background
pixels. The default is true.

\keywordindex {outline}
\keyword {outline} adds a final edge-outline pass. The default is false.

\keywordindex {outlinethreshold}
\keyword {outlinethreshold} sets the edge strength required for the outline.
The default is \type {.85}.

\keywordindex {outlineradius}
\keyword {outlineradius} sets the radius of outline dots in raster pixels. The
default is zero.

\keywordindex {luminance}
\keyword {luminance} sets the \RGB\ weights used to convert rendered colors to
intensity. The default is \type {(.2126,.7152,.0722)}.

\stopmixedcolumns

The numeric form of \type {levels} is convenient for ordinary use. A custom form
can instead provide a table of levels, each with a \type {value}, an optional
\type {radius}, and optional \type {red}, \type {green}, and \type {blue}
values. Use it when you need custom coverage values, dot sizes, or colors; use
the numeric form when evenly spaced levels are enough.

Stippling is applied to the final raster, so raster size and visible detail both
matter. Inspect the result at its intended output size and adjust the resolution,
dot radii, tone levels, and edge settings together. A stippled scene is rendered
again when needed instead of being reused from the scene cache.

\stopsection

\startsection[reference=external-meshes,title=External meshes (advanced)]

An external mesh supplies its triangles from a file instead of sampling a
function in the scene description. Add one with

\starttyping[option=MP]
lmt_scene_external [ ... ]
\stoptyping

The external loader reads standard binary STL files and the \type {lmtmesh}
representation used by \CONTEXT. It does not support ASCII STL. The file name is
passed through \CONTEXT's file resolver, so you can give a file in the job
directory or in a searched \CONTEXT\ tree by name. A path also works.

The example uses \type {luametafun-threed-sphere.stl}, a binary STL file shipped
with the \CONTEXT\ distribution. The automatic camera settings fit the imported
mesh after its bounds have been read.

\startbuffer[external-sphere]
lmt_scene_start [
  width      = 6cm,
  resolution = 120,
  crop       = true,
  projection = "perspective",
  eye        = "auto",
  target     = "auto",
  fov        = 60,
  light      = (-1,-1,-2),
  intensity  = 1.1,
] ;

lmt_scene_material [
  name      = "sphere",
  diffuse   = (.72,.82,1),
  specular  = (.7,.7,.7),
  shininess = 30,
] ;

lmt_scene_external [
  filename = "luametafun-threed-sphere.stl",
  id       = "sphere",
  material = "sphere",
] ;

lmt_scene_stop ;
\stopbuffer

\typebuffer[external-sphere][option=MP]

The external-mesh keywords are:

\startmixedcolumns

\keywordindex {filename}
\keyword {filename} names the external mesh file. The loader accepts
standard binary STL and \type {lmtmesh} files; use a path or a file name that
\CONTEXT\ can resolve.

\keywordindex {id}
\keyword {id} gives the imported object a scene name. Use a distinct name when
the object is queried later or inspected with \type {lmt_scene_mesh}.

\keywordindex {material}
\keyword {material} selects the named material used when the mesh is rendered.

\keywordindex {normal}
\keyword {normal} selects how imported normals are handled. The default smooth
mode uses supplied mesh normals when they are available; \type {flat} discards
them and shades each triangle as an independent face.

\keywordindex {flipnormals}
\keyword {flipnormals} reverses supplied packed normals before shading. Use it
when the file's orientation is opposite to the intended outside of the object.

\stopmixedcolumns

\startplacefloat
  [figure]
  [title={A binary STL mesh rendered as an ordinary shaded scene.}]
  \processMPbuffer[external-sphere]
\stopplacefloat

The coordinates in the file are used as supplied. The preset still contains a
\type {scale} key for compatibility, but the implementation does not apply it;
scale the mesh before importing if needed. Binary STL facet normals are
preserved when the file supplies nonzero normals, so their quality and
orientation affect the shading. Use \type {normal = "flat"} when those normals
are missing or unreliable, or \type {flipnormals = true} when their direction is
reversed. The imported triangles otherwise use the same camera, lighting,
material, depth handling, and cropping as other scene objects.

Large external meshes can contain many more triangles than the sampled examples
in this manual. Start with a modest output \type {resolution}, and increase it
only when the final figure needs more raster detail; raster resolution does not
reduce the cost of loading or transforming the mesh.

\stopsection

\startsection[reference=internal-meshes,title=Internal meshes (advanced)]

An internal mesh is an extension point for geometry that already exists in \LUA\
or is easier to generate there than with one of the mathematical surface
commands. A named generator returns vertices and either triangle indices or a
compact mesh type; it can also return one normal per triangle. The mesh can then
be placed, shaded, and combined with the other scene objects.

You can skip this section for ordinary graphs, parametric surfaces, and implicit
surfaces. It requires \LUA\ code and a registered mesh generator.

Here is a small pyramid generator:

\starttyping[option=LUA]
metapost.registermesher("manual-pyramid", function(specification)
    return {
        vertices = {
            { -1, -1, 0 },
            {  1, -1, 0 },
            {  1,  1, 0 },
            { -1,  1, 0 },
            {  0,  0, 1.5 },
        },
        triangles = {
            { 1, 2, 3 }, { 1, 3, 4 },
            { 1, 5, 2 }, { 2, 5, 3 },
            { 3, 5, 4 }, { 4, 5, 1 },
        },
    }
end)
\stoptyping

\startluacode
metapost.registermesher("manual-pyramid", function(specification)
    return {
        vertices = {
            { -1, -1, 0 },
            {  1, -1, 0 },
            {  1,  1, 0 },
            { -1,  1, 0 },
            {  0,  0, 1.5 },
        },
        triangles = {
            { 1, 2, 3 }, { 1, 3, 4 },
            { 1, 5, 2 }, { 2, 5, 3 },
            { 3, 5, 4 }, { 4, 5, 1 },
        },
    }
end)
\stopluacode

The indices in \type {triangles} are one-based \LUA\ indices into
\type {vertices}. The \type {triangles} field can be omitted when the generator
returns \type {meshtype}: type \type {5} treats every three consecutive vertices
as one triangle, while type \type {7} treats every four consecutive vertices as
a quadrilateral split into two triangles. Explicit indices are usually the
clearest choice for irregular meshes. An optional \type {normals} field supplies
one normal per triangle, either as a table of triplets or as a packed
\type {vector.normals} container.

After registering the generator, select it with \type {name}:

\startbuffer[internal-mesh]
lmt_scene_start [
  width = 6cm,
  crop  = true,
  eye   = (3,-4,2.5),
] ;

lmt_scene_material [
  name      = "pyramid",
  diffuse   = (.85,.35,.12),
  specular  = (.2,.2,.2),
  shininess = 20,
] ;

lmt_scene_internal [
  id       = "pyramid",
  name     = "manual-pyramid",
  normal   = "flat",
  material = "pyramid",
] ;

lmt_scene_stop ;
\stopbuffer

% \startplacefloat
%   [figure]
%   [title={An internal mesh supplied by a registered \LUA\ generator.}]
%   \processMPbuffer[internal-mesh]
% \stopplacefloat

\startplacefloat
  [figure]
  [title={A}]
  \processMPbuffer[internal-mesh]
\stopplacefloat

The internal-mesh keywords are:

\startmixedcolumns

\keywordindex {id}
\keyword {id} gives the generated mesh a name inside the scene, so that an
intersection can refer to it.

\keywordindex {name}
\keyword {name} selects the \LUA\ generator registered with
\type {metapost.registermesher}.

\keywordindex {material}
\keyword {material} selects the named material used to render the mesh.

\keywordindex {normal}
\keyword {normal} selects the normal mode for the generated mesh. In particular,
\type {flat} expands the triangles so that each face is shaded independently;
otherwise supplied triangle normals can be used for smooth shading.

\keywordindex {flipnormals}
\keyword {flipnormals} reverses a packed normal container returned by the
generator before it is used for shading.

\stopmixedcolumns

The generator receives the complete parameter specification, so it can define
settings for a particular mesh. Those settings belong to the generator, not to
the general \type {lmt_scene_internal} command.

\stopsection

\startsection[reference=lua-helpers,title=Low-level \LUA\ helpers (advanced)]

The scene commands are normally enough. \LUA\ code that generates meshes or
postprocesses a rendered image can also use two lower-level userdata libraries.
The \type {vector} library stores points, normals, and indexed meshes; the
separate \type {zbuffer} library stores a rendered pixel buffer. This section
covers the small subset needed to connect a \LUA\ mesher or image postprocessor to
the scene code. It is not a complete reference to every operation exposed by
the engine.

A registered internal mesher can return packed containers directly. This is the
packed equivalent of the table-based pyramid in the previous section:

\starttyping[option=LUA]
local vertices = vector.points.new(5,1) -- default: xyz+n
vertices(-1,-1,0)
vertices( 1,-1,0)
vertices( 1, 1,0)
vertices(-1, 1,0)
vertices( 0, 0,1.5)

local triangles = vector.mesh.new(6) -- indexed triangles
triangles(1,2,3)
triangles(1,3,4)
triangles(1,5,2)
triangles(2,5,3)
triangles(3,5,4)
triangles(4,5,1)

metapost.registermesher("packed-pyramid", function(specification)
    return {
        vertices = vertices,
        triangles = triangles,
    }
end)
\stoptyping

Use \type {name = "packed-pyramid"} with \type {lmt_scene_internal} as in the
previous section. The default point type is \type {0}, with the layout
\type {xyz+n}; it is the layout expected by the low-level triangle rasterizer.
Types \type {1} (\type {xyz}) and \type {2} (\type {xy}) are useful for
coordinate-only calculations, but should not be passed directly as the vertex
container for rasterized 3D triangles. When the optional normal components are
not supplied, they remain zero and the selected scene normal mode determines
how the mesh is handled.

\page

The main container operations are:

\startmixedcolumns

\keywordindex {vector.point}
\keyword {vector.point} creates one point from coordinates or a small \LUA\ table.
Points have one-based components, support arithmetic, and provide
\type {dotproduct}, \type {crossproduct}, \type {normalize}, \type {distance},
and \type {length}.

\keywordindex {vector.points.new}
\keyword {vector.points.new} takes \type {rows}, \type {columns}, an optional
point type, and an optional growth step. Its initial capacity is
\type {rows*columns}; call the container with coordinates to fill entries in
order. \type {get}, \type {getxyz}, \type {getnormal}, \type {set},
\type {totable}, and \type {getbounds} provide ordinary access. Indices are
one-based.

\keywordindex {vector.mesh.new}
\keyword {vector.mesh.new} takes the number of entries and an optional mesh
type. The default type is \type {3}, an indexed triangle. Call the container
with three one-based vertex indices, or use \type {get}, \type {set}, and
\type {totable}.

\keywordindex {vector.normals}
\keyword {vector.normals} creates a separate normal container with
\type {new}; its entries can be filled or changed with \type {set}, inspected
with \type {get}, and reversed with \type {flip}.

\stopmixedcolumns

For postprocessing, allocate and configure a z-buffer directly:

\starttyping[option=LUA]
local buffer = zbuffer.new()

zbuffer.setup(buffer, {
    width       = 5,
    height      = 4,
    perspective = false,
    eye         = { 0, 0, 5 },
    target      = { 0, 0, 0 },
    up          = { 0, 1, 0 },
    viewport    = { width = 5, height = 4 },
    camera      = { near = 0, far = 10, scale = 1 },
})

local visible, behind, vx, vy, vz, px, py =
    zbuffer.project(buffer, { 0, 0, 0 })
\stoptyping

The top-level \type {width} and \type {height} allocate the buffer.
\type {viewport} controls the projected pixel area, while \type {eye},
\type {target}, and \type {up} define the camera basis. With an orthographic
setup, \type {camera.scale} controls the view scale; a perspective setup uses
\type {camera.tanhalf} instead. \type {project} returns camera-space
coordinates \type {vx}, \type {vy}, and \type {vz}, followed by projected
coordinates \type {px} and \type {py}. A point behind the camera returns only
\type {false, true}.

The pixel accessors \type {get}, \type {set}, \type {getcolor},
\type {setcolor}, \type {getdepth}, \type {setdepth}, \type {getnormal}, and
\type {setnormal} use one-based \type {row,column} coordinates. \type {getsize}
returns the number of rows and columns. \type {transform} converts a point
to camera-space coordinates, and \type {tobytes} returns the \RGB\ byte string
together with its height, width, and channel count. The separate global name
\type {zbuffer} is intentional; it is not a field of \type {vector}.

These calls are enough for a controlled \LUA\ integration or a diagnostic
postprocessor. Matrix operations, contour builders, direct triangle rasterization,
and stipple stages remain implementation-facing and are not listed as a complete
low-level API.

\stopsection

\startsection[reference=raster-postprocessing,title=Raster postprocessing hooks (advanced)]

The scene \type {process} setting selects a named \LUA\ hook that runs after the
scene has been rasterized and before supersampling is resolved. It is useful for
color effects or diagnostics that need access to the rendered colors, normals,
or texture coordinates. Register the hook before using the scene:

\startluacode
local function dim(r,g,b)
    return .8 * r, .8 * g, .8 * b
end

metapost.registermeshupper {
    name    = "dim",
    color   = true,
    process = dim,
}
\stopluacode

Select it by name in the scene:

\starttyping[option=MP]
lmt_scene_start [
  process = "dim",
  ...
] ;
\stoptyping

The following rendered example uses that hook; the scene is otherwise an
ordinary surface scene. The dimmed result makes the data flow through
\type {process = "dim"} visible in the document rather than only in source.

\startbuffer[dim-scene]
lmt_scene_start [
  width      = 5.5cm,
  resolution = 120,
  crop       = true,
  eye        = (3,-4,2.6),
  process    = "dim",
] ;

lmt_scene_surface [
  code = ".7*exp(-.6*(x*x+y*y)) - .2",
  xmin = -1.8,
  xmax =  1.8,
  ymin = -1.8,
  ymax =  1.8,
  nx   = 48,
  ny   = 48,
] ;

lmt_scene_stop ;
\stopbuffer

\typebuffer[dim-scene][option=MP]

\startplacefloat
  [figure]
  [title={A scene rendered through the dim postprocessing hook.}]
  \processMPbuffer[dim-scene]
\stopplacefloat

The registration table accepts \type {name} and a required \type {process}
callback. Set \type {color}, \type {normal}, \type {texture}, or \type {slot}
to request the corresponding inputs. The callback arguments are concatenated in
that order: \type {slot} supplies zero-based \type {(row,column)}, \type {color}
supplies \type {(r,g,b)}, \type {normal} supplies \type {(nx,ny,nz)}, and
\type {texture} supplies \type {(x,y,z,t)}. For example, a callback with
\type {color = true} and \type {normal = true} receives
\type {(r,g,b,nx,ny,nz)}. It returns up to three \RGB\ values; returned numeric
values are clamped to the \type {0}--\type {1} range.

By default, background pixels are skipped. Set \type {skip = false} when the
background should also be passed to the callback. Optional \type {initialize}
and \type {finalize} functions receive the registration table before and after
the pixel pass.

For direct access to the z-buffer, set \type {direct = true} or
\type {method = "direct"}. In that form the callback receives the z-buffer and
the registration table instead of per-pixel values; use the z-buffer operations
described in the low-level \LUA\ helpers section.

The hook runs before stippling and is rerun when the scene must be rendered
again. Keep it limited to the raster work that belongs in this stage.

\stopsection

\startsection[reference=special-overlays,title=Special overlays (advanced)]

\type {lmt_scene_special} adds small, depth-tested overlays to a scene. The
supported kinds are \type {point}, \type {segment}, \type {normal}, and
\type {arrow}. Add an overlay after the object it refers to and before
\type {lmt_scene_stop}. The \type {normal} form needs a named parametric object
and its derivatives.

\startbuffer[scene-special]
lmt_scene_start [
  width      = 6cm,
  resolution = 120,
  crop       = true,
  eye        = (3,-4,2.6),
] ;

lmt_scene_parametric [
  id  = "surface",
  x   = "u",
  y   = "v",
  z   = ".7*exp(-.6*(u*u+v*v)) - .2",
  xu  = "1",
  yu  = "0",
  zu  = "-.84*u*exp(-.6*(u*u+v*v))",
  xv  = "0",
  yv  = "1",
  zv  = "-.84*v*exp(-.6*(u*u+v*v))",
  umin = -1.8,
  umax =  1.8,
  vmin = -1.8,
  vmax =  1.8,
  nu   = 48,
  nv   = 48,
] ;

lmt_scene_special [
  kind  = "point",
  id    = "probe",
  at    = (0,0,.5),
  color = red,
  width = 2,
] ;

lmt_scene_special [
  kind  = "segment",
  id    = "baseline",
  from  = (-1.2,-1,0),
  to    = ( 1.2,-1,0),
  color = blue,
  width = 1,
] ;

lmt_scene_special [
  kind   = "normal",
  id     = "normal",
  name   = "surface",
  u      = 0,
  v      = 0,
  length = .2,
  span   = .12,
  color  = green,
  width  = 1,
] ;

lmt_scene_special [
  kind   = "arrow",
  id     = "arrow",
  from   = (0,0,0),
  to     = (0,0,1),
  length = .15,
  span   = .12,
  color  = (.9,.4,0),
  width  = 1,
] ;

lmt_scene_stop ;
\stopbuffer

\typebuffer[scene-special][option=MP]

\startplacefloat
  [figure]
  [title={Point, segment, normal, and arrow overlays in one parametric scene.}]
  \processMPbuffer[scene-special]
\stopplacefloat

The special-overlay keywords are:

\startmixedcolumns

\keywordindex {kind}
\keyword {kind} selects the overlay form: \type {point}, \type {segment},
\type {normal}, or \type {arrow}.

\keywordindex {id}
\keyword {id} gives the overlay a name.

\keywordindex {at}
\keyword {at} gives the literal 3D point where a \type {point} overlay is
drawn. A point overlay draws a small cross there.

\keywordindex {point}
\keywordindex {x}
\keywordindex {y}
\keywordindex {z}
\keyword {point} gives a spatial coordinate when a named object is resolved.
The same coordinate can be given with \type {x}, \type {y}, and \type {z}.

\keywordindex {from}
\keyword {from} gives the first endpoint of a \type {segment} or \type {arrow}.

\keywordindex {to}
\keyword {to} gives the second endpoint of a \type {segment} or \type {arrow}.

\keywordindex {name}
\keyword {name} selects a named scene object for point or normal lookup. It is
required by the \type {normal} form.

\keywordindex {u}
\keyword {u} gives the first parameter when a named parametric object is
resolved.

\keywordindex {v}
\keyword {v} gives the second parameter when a named parametric object is
resolved.

\keywordindex {t}
\keyword {t} gives the parameter when a named curve is resolved.

\keywordindex {length}
\keyword {length} controls the size of a point cross or the length of an
arrowhead.

\keywordindex {span}
\keyword {span} controls the width of an arrowhead.

\keywordindex {rotation}
\keyword {rotation} rotates an arrowhead around its direction.

\keywordindex {color}
\keyword {color} selects the overlay line color.

\keywordindex {width}
\keyword {width} sets the overlay line width in raster pixels.

\stopmixedcolumns

The \type {normal} form resolves the named object using the same coordinate or
parameter lookup and evaluates its normal; supplying analytic derivatives makes
a parametric normal reliable. All four forms are depth-tested like other scene
geometry.

The named-normal example uses a parametric surface. For a graph or implicit
surface, use the direct \type {lmt_scene_normal} query after rendering.

\stopsection

\startsection[reference=mesh-inspection,title=Mesh inspection (experimental)]

This final example is diagnostic only. It shows the triangle paths captured for
a scene object, which helps you check sampling, winding, and mesh coverage. It
is not a normal rendering interface. Set the object's experimental \type {vector}
flag, render the scene, and draw \type {lmt_scene_mesh} using its \type {id}:

\startbuffer[mesh-inspection]
lmt_scene_start [
  width      = 6cm,
  resolution = 120,
  crop       = true,
  eye        = (3,-4,2.6),
] ;

lmt_scene_surface [
  id       = "surface",
  vector   = true,
  code     = ".7*exp(-.6*(x*x+y*y)) - .2",
  dfdx     = "-.84*x*exp(-.6*(x*x+y*y))",
  dfdy     = "-.84*y*exp(-.6*(x*x+y*y))",
  xmin     = -1.8,
  xmax     =  1.8,
  ymin     = -1.8,
  ymax     =  1.8,
  nx       = 24,
  ny       = 24,
] ;

lmt_scene_render ;

draw lmt_scene_mesh "surface"
  withpen pencircle scaled .2pt
  withcolor red ;

lmt_scene_stop ;
\stopbuffer

\typebuffer[mesh-inspection][option=MP]

\startplacefloat
  [figure]
  [title={Diagnostic view of the sampled triangle mesh.}]
  \processMPbuffer[mesh-inspection]
\stopplacefloat

\type {lmt_scene_mesh} returns a path made from the captured projected triangles.
The path can include triangles hidden in the shaded rendering, so use it to
inspect the mesh rather than to produce a clean visible outline. Call it after
\type {lmt_scene_render}; the object's \type {vector = true} flag captures the
paths. This flag and helper depend on experimental vector-point support and may
change, so use them while debugging a scene.

\stopsection

\startsection[reference=rendering-pipeline,title=Rendering pipeline (advanced)]

\index {mesh counts}%
\index {raster resolution}%
\index {depth testing}%
The scene commands are declarative, but a rendered figure passes through several
distinct stages. This distinction is useful when deciding whether to increase
an object's mesh counts, the raster resolution, or the supersampling factor. It
also explains why rendering must happen before projection and normal queries can
be used.

The main sources of render time are dense implicit grids
(\type {nx*ny*nz}), large sampled surface grids, high \type {supersample} values
which enlarge the temporary raster area roughly with the square of the factor,
large external meshes, and mesh--mesh intersections. When a figure is slow,
inspect these sources separately and increase one of them at a time.

For an ordinary render without a cache hit, the overall path is:

\starttyping
scene description
    -> object generation and triangle meshes
    -> camera, crop, and raster setup
    -> projection and z-buffer rasterization
    -> resolution and raster postprocessing
    -> final crop and bytemap conversion
    -> MetaPost placement and scene queries
\stoptyping

The stages have the following roles:

\startitemize
  \startitem
    The scene description records the shared camera, lighting, materials, output
    settings, and the objects that belong to the scene. Object commands do not
    immediately draw pixels; they add specifications to the scene.
  \stopitem
  \startitem
    Object generation turns those specifications into geometry. A graph or
    parametric surface is sampled over its parameter grid, an implicit surface is
    extracted from samples in a 3D box, and external or internal mesh objects
    provide their triangles directly. Normals are generated or attached at this
    stage so that they are available during shading. The object mesh counts
    control this geometric stage.
  \stopitem
  \startitem
    Camera and raster setup chooses the internal pixel dimensions, allocates the
    z-buffer, and derives the camera transform. \type {width}, \type {height},
    \type {bytewidth}, \type {byteheight}, and \type {resolution} determine the
    displayed and raster dimensions; \type {supersample} enlarges the temporary
    raster used while drawing. With automatic cropping, the projected object
    bounds also help fit the camera transform to the scene.
  \stopitem
  \startitem
    Projection transforms mesh vertices into camera and screen coordinates.
    Triangle rasterization then writes color, depth, and normal information to
    the z-buffer. Depth tests decide which surfaces are visible, while materials,
    lights, opacity, and the selected normals determine the shaded result.
    Transparent and on-top objects are handled in their corresponding rendering
    passes.
  \stopitem
  \startitem
    After the internal drawing pass, an optional \type {process} hook (described
    in the \goto {Raster postprocessing hooks} [raster-postprocessing] section)
    can inspect or modify the rendered buffer. The supersampled buffer is then
    reduced to the display raster. Stippling, if enabled, is applied later to
    that resolved raster; neither operation changes the generated triangle mesh.
  \stopitem
  \startitem
    Cropping then trims the resolved raster to the occupied scene bounds when
    requested. Thus \type {crop} affects both the camera fitting used before
    rasterization and, where appropriate, the dimensions of the final displayed
    buffer. The resulting pixel data is converted to a \METAPOST\ bytemap.
  \stopitem
  \startitem
    \type {lmt_scene_render} makes the rendered result available to projection,
    visibility, bounds, bytemap, and normal helpers. \type {lmt_scene_stop}
    performs the render when needed and places the finished bytemap in the
    current picture. Queries and ordinary \METAPOST\ annotations therefore come
    after rendering, while the scene is still open.
  \stopitem
\stopitemize

This separation also explains several common adjustments. Increasing
\type {nx}, \type {ny}, \type {nz}, \type {nu}, or \type {nv} changes the
geometric approximation and can reveal features that a coarse mesh misses.
Increasing \type {resolution} or \type {supersample} changes the raster quality
of an already generated scene; it does not recover geometric detail that was
never sampled. Stippling similarly changes the final raster rather than the
surface or its projection.

When caching is enabled and a valid cached result is available, the saved
rendered scene can be restored instead of repeating the ordinary generation and
rasterization stages. The cache is scene-level: it is not a reusable cache of
individual object meshes.

\stopsection

\startsection[title=About the cover page]

The cover captures the rendered scene in a \METAPOST\ picture with
\type {picture q := image(...)}. The scene is rendered while the argument of
\type {image} is evaluated; afterwards, shifting by \type {-center q} moves the
picture's center to the origin, and the final shift places that center on the
paper. This is a useful technique when a scene must be captured, scaled, and
positioned independently of the surrounding \METAPOST\ drawing.

\typebuffer[coverpage][option=TEX]

\stopsection

\starttitle[title=Index]

\placeregister[index,keywordindex]

\stoptitle

\stoptext
